%%%%%%%%%%%%%%%%%%%%%%%%%%%%%%%%%%%%%%%%%%%%%%%%%%%%%%%%%%%%
% 11.11 HIGH-PERFORMANCE MATHEMATICAL COMPUTING
% The Composition of Advanced Mathematics
%%%%%%%%%%%%%%%%%%%%%%%%%%%%%%%%%%%%%%%%%%%%%%%%%%%%%%%%%%%%

\section{High-Performance Mathematical Computing}
\label{sec:high-performance-mathematical-computing}

\prerequisite{
Scientific computing, numerical linear algebra, numerical differential
equations, Monte Carlo methods, sparse matrices, algorithmic complexity,
parallel computing concepts, and basic computer architecture.
}

\learningobjectives{
By the end of this section, the reader should be able to:
\begin{itemize}
    \item explain the purpose of high-performance mathematical computing;
    \item distinguish algorithmic, implementation, and hardware
          performance;
    \item understand computational throughput;
    \item recognize floating-point operation counts;
    \item understand memory bandwidth limitations;
    \item interpret arithmetic intensity;
    \item understand the roofline model conceptually;
    \item distinguish latency-bound, bandwidth-bound, and compute-bound
          algorithms;
    \item understand cache blocking and tiling;
    \item recognize vectorization and SIMD computation;
    \item understand dense linear algebra performance;
    \item recognize matrix-matrix multiplication as a high-intensity
          computational kernel;
    \item understand sparse linear algebra performance;
    \item recognize irregular memory access;
    \item understand matrix-free computation;
    \item understand shared-memory parallelism;
    \item understand distributed-memory parallelism;
    \item recognize message-passing patterns;
    \item distinguish local and global communication;
    \item understand collective operations;
    \item recognize reductions and broadcasts;
    \item understand communication complexity;
    \item analyze domain decomposition;
    \item understand halo exchange;
    \item recognize strong and weak scaling;
    \item calculate speedup and parallel efficiency;
    \item apply Amdahl's law;
    \item understand Gustafson-type scaling conceptually;
    \item recognize load imbalance;
    \item understand synchronization costs;
    \item recognize communication hiding;
    \item understand asynchronous computation conceptually;
    \item recognize accelerator computing;
    \item understand GPU execution conceptually;
    \item distinguish host and accelerator memory;
    \item understand data-transfer overhead;
    \item recognize kernel launch overhead;
    \item understand batch and fused operations;
    \item recognize mixed-precision computation;
    \item understand iterative refinement;
    \item recognize tensor and accelerator arithmetic;
    \item understand parallel numerical linear algebra;
    \item recognize distributed matrix layouts;
    \item understand block decomposition;
    \item recognize communication-avoiding algorithms;
    \item understand scalable Krylov methods;
    \item recognize global-reduction bottlenecks;
    \item understand multigrid scalability;
    \item recognize parallel FFT algorithms;
    \item understand distributed PDE solvers;
    \item recognize parallel Monte Carlo;
    \item understand ensemble parallelism;
    \item recognize parallel optimization;
    \item understand checkpointing at scale;
    \item recognize resilience and fault tolerance;
    \item understand reproducibility in parallel floating-point
          computation;
    \item distinguish bitwise and numerical reproducibility;
    \item understand energy and resource efficiency conceptually;
    \item recognize performance portability;
    \item understand benchmarking and profiling at scale;
    \item design scalable computational experiments;
    \item connect high-performance computing with simulation,
          optimization, numerical linear algebra, PDEs, statistics,
          probability, and applied mathematics.
\end{itemize}
}

%%%%%%%%%%%%%%%%%%%%%%%%%%%%%%%%%%%%%%%%%%%%%%%%%%%%%%%%%%%%
\subsection{What Is High-Performance Mathematical Computing?}
\label{subsec:what-is-hpmc}
%%%%%%%%%%%%%%%%%%%%%%%%%%%%%%%%%%%%%%%%%%%%%%%%%%%%%%%%%%%%

High-performance mathematical computing studies how mathematical
algorithms can be executed efficiently on large and parallel computer
systems.

The central goal is not merely

\[
\boxed{
\text{to compute quickly},
}
\]

but rather

\[
\boxed{
\text{to solve mathematically meaningful problems at scales that would
otherwise be impractical}.
}
\]

This includes problems involving:

\[
\boxed{
\text{large matrices},
}
\]

\[
\boxed{
\text{fine computational meshes},
}
\]

\[
\boxed{
\text{large simulation ensembles},
}
\]

and

\[
\boxed{
\text{high-dimensional optimization}.
}
\]

%%%%%%%%%%%%%%%%%%%%%%%%%%%%%%%%%%%%%%%%%%%%%%%%%%%%%%%%%%%%
\subsection{Performance Has Several Layers}
\label{subsec:hpmc-performance-layers}
%%%%%%%%%%%%%%%%%%%%%%%%%%%%%%%%%%%%%%%%%%%%%%%%%%%%%%%%%%%%

Computational performance depends on several interacting layers:

\[
\boxed{
\text{mathematical formulation}
}
\]

\[
\downarrow
\]

\[
\boxed{
\text{numerical algorithm}
}
\]

\[
\downarrow
\]

\[
\boxed{
\text{software implementation}
}
\]

\[
\downarrow
\]

\[
\boxed{
\text{hardware architecture}.
}
\]

Improving one layer cannot always compensate for a poor choice at another
layer.

%%%%%%%%%%%%%%%%%%%%%%%%%%%%%%%%%%%%%%%%%%%%%%%%%%%%%%%%%%%%
\subsection{Algorithmic Improvement Versus Faster Hardware}
\label{subsec:algorithm-versus-hardware}
%%%%%%%%%%%%%%%%%%%%%%%%%%%%%%%%%%%%%%%%%%%%%%%%%%%%%%%%%%%%

Suppose one algorithm costs

\[
O(n^3)
\]

while another costs

\[
O(n^2).
\]

For sufficiently large \(n\), the algorithmic improvement dominates any
fixed hardware speedup.

Thus

\[
\boxed{
\text{better mathematics}
}
\]

can provide a larger performance gain than

\[
\boxed{
\text{faster processors}.
}
\]

High-performance computing begins with good algorithmic design.

%%%%%%%%%%%%%%%%%%%%%%%%%%%%%%%%%%%%%%%%%%%%%%%%%%%%%%%%%%%%
\subsection{Floating-Point Operations}
\label{subsec:hpmc-flops}
%%%%%%%%%%%%%%%%%%%%%%%%%%%%%%%%%%%%%%%%%%%%%%%%%%%%%%%%%%%%

A floating-point operation, abbreviated FLOP, is an arithmetic operation
performed on floating-point values.

Performance is often expressed as

\[
\boxed{
\text{FLOP/s}
=
\frac{
\text{number of floating-point operations}
}{
\text{execution time}
}.
}
\]

Large values indicate high arithmetic throughput.

However, FLOP/s alone does not characterize every algorithm.

%%%%%%%%%%%%%%%%%%%%%%%%%%%%%%%%%%%%%%%%%%%%%%%%%%%%%%%%%%%%
\subsection{Peak and Sustained Performance}
\label{subsec:peak-sustained-performance}
%%%%%%%%%%%%%%%%%%%%%%%%%%%%%%%%%%%%%%%%%%%%%%%%%%%%%%%%%%%%

Peak performance is the theoretical maximum arithmetic rate of the
hardware.

Sustained performance is the rate actually achieved by an application.

Usually,

\[
\boxed{
\text{sustained performance}
<
\text{peak performance}.
}
\]

The gap may arise from:

\[
\boxed{
\text{memory traffic},
}
\]

\[
\boxed{
\text{communication},
}
\]

\[
\boxed{
\text{branching},
}
\]

or

\[
\boxed{
\text{poor parallel utilization}.
}
\]

%%%%%%%%%%%%%%%%%%%%%%%%%%%%%%%%%%%%%%%%%%%%%%%%%%%%%%%%%%%%
\subsection{Memory Bandwidth}
\label{subsec:hpmc-memory-bandwidth}
%%%%%%%%%%%%%%%%%%%%%%%%%%%%%%%%%%%%%%%%%%%%%%%%%%%%%%%%%%%%

Memory bandwidth measures the rate at which data can be transferred
between memory and processing units.

A simple model is

\[
\boxed{
\text{time}
\approx
\frac{
\text{bytes transferred}
}{
\text{memory bandwidth}
}.
}
\]

Algorithms performing little arithmetic per byte may be limited primarily
by memory bandwidth rather than arithmetic capability.

%%%%%%%%%%%%%%%%%%%%%%%%%%%%%%%%%%%%%%%%%%%%%%%%%%%%%%%%%%%%
\subsection{Arithmetic Intensity}
\label{subsec:hpmc-arithmetic-intensity}
%%%%%%%%%%%%%%%%%%%%%%%%%%%%%%%%%%%%%%%%%%%%%%%%%%%%%%%%%%%%

Arithmetic intensity is

\[
\boxed{
I
=
\frac{
\text{floating-point operations}
}{
\text{bytes transferred}
}.
}
\]

Low \(I\) suggests a memory-bound computation.

High \(I\) provides greater opportunity to use the processor's arithmetic
capability.

%%%%%%%%%%%%%%%%%%%%%%%%%%%%%%%%%%%%%%%%%%%%%%%%%%%%%%%%%%%%
\subsection{Roofline Model}
\label{subsec:roofline-model}
%%%%%%%%%%%%%%%%%%%%%%%%%%%%%%%%%%%%%%%%%%%%%%%%%%%%%%%%%%%%

A simplified roofline performance model estimates attainable performance
as

\[
\boxed{
P
\leq
\min
\left(
P_{\mathrm{peak}},
B I
\right),
}
\]

where

\[
P_{\mathrm{peak}}
\]

is peak arithmetic throughput,

\[
B
\]

is memory bandwidth, and

\[
I
\]

is arithmetic intensity.

The model separates two main regimes.

%%%%%%%%%%%%%%%%%%%%%%%%%%%%%%%%%%%%%%%%%%%%%%%%%%%%%%%%%%%%
\subsection{Bandwidth-Bound Regime}
\label{subsec:bandwidth-bound-regime}
%%%%%%%%%%%%%%%%%%%%%%%%%%%%%%%%%%%%%%%%%%%%%%%%%%%%%%%%%%%%

If

\[
BI<P_{\mathrm{peak}},
\]

then

\[
\boxed{
P
\lesssim
BI.
}
\]

Performance is limited primarily by memory bandwidth.

Increasing arithmetic capability alone produces little benefit.

%%%%%%%%%%%%%%%%%%%%%%%%%%%%%%%%%%%%%%%%%%%%%%%%%%%%%%%%%%%%
\subsection{Compute-Bound Regime}
\label{subsec:compute-bound-regime}
%%%%%%%%%%%%%%%%%%%%%%%%%%%%%%%%%%%%%%%%%%%%%%%%%%%%%%%%%%%%

If

\[
BI\geq P_{\mathrm{peak}},
\]

then

\[
\boxed{
P
\lesssim
P_{\mathrm{peak}}.
}
\]

The computation has enough arithmetic intensity to approach the hardware's
computational ceiling.

%%%%%%%%%%%%%%%%%%%%%%%%%%%%%%%%%%%%%%%%%%%%%%%%%%%%%%%%%%%%
\subsection{Worked Example: Roofline Estimate}
\label{subsec:worked-roofline-estimate}
%%%%%%%%%%%%%%%%%%%%%%%%%%%%%%%%%%%%%%%%%%%%%%%%%%%%%%%%%%%%

\begin{workedexamplebox}{Estimating a Performance Bound}

Suppose a processor has peak performance

\[
P_{\mathrm{peak}}
=
1000
\]

GFLOP/s and memory bandwidth

\[
B
=
200
\]

GB/s.

Suppose an algorithm has arithmetic intensity

\[
I
=
2
\]

FLOP/byte.

Then

\[
BI
=
200(2)
=
400
\]

GFLOP/s.

Therefore,

\[
P
\leq
\min(1000,400).
\]

\begin{answerbox}
\[
\boxed{
P
\leq
400\text{ GFLOP/s}.
}
\]
\end{answerbox}

The algorithm is bandwidth bound under this simplified model.

\end{workedexamplebox}

%%%%%%%%%%%%%%%%%%%%%%%%%%%%%%%%%%%%%%%%%%%%%%%%%%%%%%%%%%%%
\subsection{Operational Intensity Threshold}
\label{subsec:roofline-threshold}
%%%%%%%%%%%%%%%%%%%%%%%%%%%%%%%%%%%%%%%%%%%%%%%%%%%%%%%%%%%%

The crossover between bandwidth-bound and compute-bound behavior occurs
near

\[
\boxed{
I^*
=
\frac{
P_{\mathrm{peak}}
}{
B
}.
}
\]

If

\[
I<I^*,
\]

the computation tends to be bandwidth bound.

If

\[
I>I^*,
\]

the computation may become compute bound.

%%%%%%%%%%%%%%%%%%%%%%%%%%%%%%%%%%%%%%%%%%%%%%%%%%%%%%%%%%%%
\subsection{Data Reuse}
\label{subsec:hpmc-data-reuse}
%%%%%%%%%%%%%%%%%%%%%%%%%%%%%%%%%%%%%%%%%%%%%%%%%%%%%%%%%%%%

Arithmetic intensity can often be improved by reusing data already loaded
into fast memory.

For example, a matrix block may be loaded once and used for many
multiplications.

Thus

\[
\boxed{
\text{data reuse}
\longrightarrow
\text{higher arithmetic intensity}.
}
\]

%%%%%%%%%%%%%%%%%%%%%%%%%%%%%%%%%%%%%%%%%%%%%%%%%%%%%%%%%%%%
\subsection{Cache Blocking}
\label{subsec:cache-blocking}
%%%%%%%%%%%%%%%%%%%%%%%%%%%%%%%%%%%%%%%%%%%%%%%%%%%%%%%%%%%%

Cache blocking divides a computation into smaller pieces that fit into
faster memory.

For matrix multiplication,

\[
C=AB,
\]

partition matrices into blocks:

\[
\boxed{
A_{ik},
\qquad
B_{kj},
\qquad
C_{ij}.
}
\]

Then update

\[
\boxed{
C_{ij}
\leftarrow
C_{ij}
+
A_{ik}B_{kj}.
}
\]

The goal is to reuse blocks before they leave cache.

%%%%%%%%%%%%%%%%%%%%%%%%%%%%%%%%%%%%%%%%%%%%%%%%%%%%%%%%%%%%
\subsection{Tiling}
\label{subsec:hpmc-tiling}
%%%%%%%%%%%%%%%%%%%%%%%%%%%%%%%%%%%%%%%%%%%%%%%%%%%%%%%%%%%%

Tiling is the general strategy of dividing iteration spaces or arrays into
blocks.

It improves locality by keeping working data close to the processor.

Tiling appears in:

\[
\boxed{
\text{dense matrix multiplication},
}
\]

\[
\boxed{
\text{stencil computation},
}
\]

and

\[
\boxed{
\text{tensor operations}.
}
\]

%%%%%%%%%%%%%%%%%%%%%%%%%%%%%%%%%%%%%%%%%%%%%%%%%%%%%%%%%%%%
\subsection{Vectorization and SIMD}
\label{subsec:vectorization-simd}
%%%%%%%%%%%%%%%%%%%%%%%%%%%%%%%%%%%%%%%%%%%%%%%%%%%%%%%%%%%%

SIMD means Single Instruction, Multiple Data.

A single processor instruction operates on several numbers simultaneously.

For vectors

\[
x,y\in\mathbb{R}^n,
\]

operations such as

\[
\boxed{
y
\leftarrow
ax+y
}
\]

can process several vector entries in parallel.

Algorithms with regular arithmetic patterns vectorize especially well.

%%%%%%%%%%%%%%%%%%%%%%%%%%%%%%%%%%%%%%%%%%%%%%%%%%%%%%%%%%%%
\subsection{Data Alignment and Contiguous Access}
\label{subsec:contiguous-memory-access}
%%%%%%%%%%%%%%%%%%%%%%%%%%%%%%%%%%%%%%%%%%%%%%%%%%%%%%%%%%%%

Processors typically move memory in blocks.

Sequential access to contiguous elements therefore tends to be efficient.

Irregular access such as

\[
x_{p_1},
x_{p_2},
x_{p_3},
\ldots
\]

with unpredictable indices can reduce locality and effective bandwidth.

This distinction is important in sparse computation.

%%%%%%%%%%%%%%%%%%%%%%%%%%%%%%%%%%%%%%%%%%%%%%%%%%%%%%%%%%%%
\subsection{Dense Matrix Multiplication}
\label{subsec:hpmc-dense-matrix-multiplication}
%%%%%%%%%%%%%%%%%%%%%%%%%%%%%%%%%%%%%%%%%%%%%%%%%%%%%%%%%%%%

For dense matrices

\[
A,B\in\mathbb{R}^{n\times n},
\]

standard matrix multiplication requires approximately

\[
\boxed{
2n^3
}
\]

floating-point operations.

The result contains only

\[
n^2
\]

entries.

Because matrix entries can be reused many times, dense matrix
multiplication can achieve high arithmetic intensity.

%%%%%%%%%%%%%%%%%%%%%%%%%%%%%%%%%%%%%%%%%%%%%%%%%%%%%%%%%%%%
\subsection{Why Matrix-Matrix Multiplication Is Efficient}
\label{subsec:why-gemm-efficient}
%%%%%%%%%%%%%%%%%%%%%%%%%%%%%%%%%%%%%%%%%%%%%%%%%%%%%%%%%%%%

In

\[
C=AB,
\]

each element of \(A\) and \(B\) participates in many arithmetic operations.

Well-blocked implementations exploit this reuse.

Therefore dense matrix multiplication is a fundamental high-performance
kernel and often approaches a large fraction of machine peak throughput.

%%%%%%%%%%%%%%%%%%%%%%%%%%%%%%%%%%%%%%%%%%%%%%%%%%%%%%%%%%%%
\subsection{Matrix-Vector Multiplication}
\label{subsec:hpmc-matrix-vector}
%%%%%%%%%%%%%%%%%%%%%%%%%%%%%%%%%%%%%%%%%%%%%%%%%%%%%%%%%%%%

For dense

\[
y=Ax,
\]

the operation count is approximately

\[
\boxed{
2n^2,
}
\]

while the entire matrix \(A\) must usually be read.

Each matrix entry contributes to only a small number of operations.

Thus matrix-vector multiplication generally has much lower arithmetic
intensity than matrix-matrix multiplication.

%%%%%%%%%%%%%%%%%%%%%%%%%%%%%%%%%%%%%%%%%%%%%%%%%%%%%%%%%%%%
\subsection{Sparse Matrix-Vector Multiplication}
\label{subsec:hpmc-spmv}
%%%%%%%%%%%%%%%%%%%%%%%%%%%%%%%%%%%%%%%%%%%%%%%%%%%%%%%%%%%%

For sparse matrix \(A\),

\[
\boxed{
y=Ax
}
\]

requires approximately

\[
\boxed{
2\operatorname{nnz}(A)
}
\]

floating-point operations.

However, the computation must also read:

\[
\boxed{
\text{values},
}
\]

\[
\boxed{
\text{column indices},
}
\]

and irregular elements of

\[
\boxed{
x.
}
\]

Sparse matrix-vector multiplication is therefore commonly bandwidth
limited.

%%%%%%%%%%%%%%%%%%%%%%%%%%%%%%%%%%%%%%%%%%%%%%%%%%%%%%%%%%%%
\subsection{Sparse Storage Formats}
\label{subsec:hpmc-sparse-storage}
%%%%%%%%%%%%%%%%%%%%%%%%%%%%%%%%%%%%%%%%%%%%%%%%%%%%%%%%%%%%

Sparse matrices may be stored using structures such as:

\[
\boxed{
\text{compressed rows},
}
\]

\[
\boxed{
\text{compressed columns},
}
\]

or

\[
\boxed{
\text{coordinate lists}.
}
\]

The best format depends on the computational operation and matrix
structure.

%%%%%%%%%%%%%%%%%%%%%%%%%%%%%%%%%%%%%%%%%%%%%%%%%%%%%%%%%%%%
\subsection{Structured Sparse Matrices}
\label{subsec:structured-sparse-matrices}
%%%%%%%%%%%%%%%%%%%%%%%%%%%%%%%%%%%%%%%%%%%%%%%%%%%%%%%%%%%%

Matrices arising from regular PDE grids often have predictable sparsity.

For example, a five-point stencil has only a few neighbors per grid point.

In such cases, it may be unnecessary to store matrix indices explicitly.

The stencil itself defines the matrix action.

%%%%%%%%%%%%%%%%%%%%%%%%%%%%%%%%%%%%%%%%%%%%%%%%%%%%%%%%%%%%
\subsection{Matrix-Free Computation}
\label{subsec:hpmc-matrix-free}
%%%%%%%%%%%%%%%%%%%%%%%%%%%%%%%%%%%%%%%%%%%%%%%%%%%%%%%%%%%%

A matrix-free method computes

\[
\boxed{
v
\mapsto
Av
}
\]

without explicitly storing \(A\).

This is particularly useful when \(A\) arises from a differential
operator.

Advantages include:

\[
\boxed{
\text{lower memory use},
}
\]

and sometimes

\[
\boxed{
\text{greater arithmetic intensity}.
}
\]

%%%%%%%%%%%%%%%%%%%%%%%%%%%%%%%%%%%%%%%%%%%%%%%%%%%%%%%%%%%%
\subsection{Matrix-Free PDE Operator}
\label{subsec:matrix-free-pde-operator}
%%%%%%%%%%%%%%%%%%%%%%%%%%%%%%%%%%%%%%%%%%%%%%%%%%%%%%%%%%%%

For the discrete one-dimensional Laplacian,

\[
(Au)_i
=
\frac{
-u_{i-1}
+
2u_i
-
u_{i+1}
}{
h^2
}.
\]

Instead of storing the tridiagonal matrix, compute this stencil directly.

The matrix is represented implicitly by the numerical operator.

%%%%%%%%%%%%%%%%%%%%%%%%%%%%%%%%%%%%%%%%%%%%%%%%%%%%%%%%%%%%
\subsection{Shared-Memory Parallelism}
\label{subsec:hpmc-shared-memory}
%%%%%%%%%%%%%%%%%%%%%%%%%%%%%%%%%%%%%%%%%%%%%%%%%%%%%%%%%%%%

In shared-memory computation, many processor cores access the same memory
space.

A loop may be divided across threads:

\[
\boxed{
\{1,\ldots,N\}
=
I_1
\cup
\cdots
\cup
I_p.
}
\]

Thread \(j\) processes indices in \(I_j\).

%%%%%%%%%%%%%%%%%%%%%%%%%%%%%%%%%%%%%%%%%%%%%%%%%%%%%%%%%%%%
\subsection{Parallel Loop Independence}
\label{subsec:parallel-loop-independence}
%%%%%%%%%%%%%%%%%%%%%%%%%%%%%%%%%%%%%%%%%%%%%%%%%%%%%%%%%%%%

A loop is easy to parallelize when iterations are independent.

For example,

\[
\boxed{
y_i
=
a x_i+b_i
}
\]

can be computed independently for every \(i\).

If iterations depend on previous results, additional algorithmic
restructuring may be required.

%%%%%%%%%%%%%%%%%%%%%%%%%%%%%%%%%%%%%%%%%%%%%%%%%%%%%%%%%%%%
\subsection{Race Conditions}
\label{subsec:hpmc-race-conditions}
%%%%%%%%%%%%%%%%%%%%%%%%%%%%%%%%%%%%%%%%%%%%%%%%%%%%%%%%%%%%

Suppose several threads execute

\[
s
\leftarrow
s+x_i.
\]

If updates occur simultaneously without coordination, increments may be
lost.

This is a race condition.

Parallel reductions require specialized synchronization or reduction
strategies.

%%%%%%%%%%%%%%%%%%%%%%%%%%%%%%%%%%%%%%%%%%%%%%%%%%%%%%%%%%%%
\subsection{Parallel Reductions}
\label{subsec:hpmc-parallel-reductions}
%%%%%%%%%%%%%%%%%%%%%%%%%%%%%%%%%%%%%%%%%%%%%%%%%%%%%%%%%%%%

To compute

\[
s
=
\sum_{i=1}^{N}
x_i,
\]

each worker can first compute a local sum

\[
\boxed{
s_j
=
\sum_{i\in I_j}
x_i.
}
\]

Then combine the local sums:

\[
\boxed{
s
=
\sum_{j=1}^{p}
s_j.
}
\]

A tree reduction reduces the parallel depth from \(O(p)\) to approximately

\[
\boxed{
O(\log p).
}
\]

%%%%%%%%%%%%%%%%%%%%%%%%%%%%%%%%%%%%%%%%%%%%%%%%%%%%%%%%%%%%
\subsection{Synchronization Overhead}
\label{subsec:hpmc-synchronization-overhead}
%%%%%%%%%%%%%%%%%%%%%%%%%%%%%%%%%%%%%%%%%%%%%%%%%%%%%%%%%%%%

Parallel tasks sometimes must wait until all participants reach the same
stage.

A synchronization barrier enforces such coordination.

Frequent barriers can reduce performance because fast workers wait for
slow workers.

Thus scalable algorithms minimize unnecessary synchronization.

%%%%%%%%%%%%%%%%%%%%%%%%%%%%%%%%%%%%%%%%%%%%%%%%%%%%%%%%%%%%
\subsection{Distributed-Memory Parallelism}
\label{subsec:hpmc-distributed-memory}
%%%%%%%%%%%%%%%%%%%%%%%%%%%%%%%%%%%%%%%%%%%%%%%%%%%%%%%%%%%%

In distributed-memory systems, each process owns separate local memory.

Data must be exchanged explicitly.

The computation therefore contains:

\[
\boxed{
\text{local arithmetic}
+
\text{communication}.
}
\]

Large clusters and supercomputers commonly use this model.

%%%%%%%%%%%%%%%%%%%%%%%%%%%%%%%%%%%%%%%%%%%%%%%%%%%%%%%%%%%%
\subsection{Point-to-Point Communication}
\label{subsec:hpmc-point-to-point}
%%%%%%%%%%%%%%%%%%%%%%%%%%%%%%%%%%%%%%%%%%%%%%%%%%%%%%%%%%%%

A process can send data directly to another process:

\[
\boxed{
P_i
\rightarrow
P_j.
}
\]

This is useful for exchanging local neighboring information.

Domain-decomposed PDE methods commonly use point-to-point communication.

%%%%%%%%%%%%%%%%%%%%%%%%%%%%%%%%%%%%%%%%%%%%%%%%%%%%%%%%%%%%
\subsection{Collective Communication}
\label{subsec:hpmc-collectives}
%%%%%%%%%%%%%%%%%%%%%%%%%%%%%%%%%%%%%%%%%%%%%%%%%%%%%%%%%%%%

Collective communication involves groups of processes.

Important examples include:

\[
\boxed{
\text{broadcast},
}
\]

\[
\boxed{
\text{reduction},
}
\]

\[
\boxed{
\text{all-reduction},
}
\]

and

\[
\boxed{
\text{gather/scatter}.
}
\]

Collectives simplify communication patterns but may become scalability
bottlenecks.

%%%%%%%%%%%%%%%%%%%%%%%%%%%%%%%%%%%%%%%%%%%%%%%%%%%%%%%%%%%%
\subsection{Broadcast}
\label{subsec:hpmc-broadcast}
%%%%%%%%%%%%%%%%%%%%%%%%%%%%%%%%%%%%%%%%%%%%%%%%%%%%%%%%%%%%

A broadcast sends one process's data to all others.

Conceptually,

\[
\boxed{
P_0
\longrightarrow
P_1,\ldots,P_{p-1}.
}
\]

Broadcasts occur when every worker requires shared parameters or matrix
blocks.

%%%%%%%%%%%%%%%%%%%%%%%%%%%%%%%%%%%%%%%%%%%%%%%%%%%%%%%%%%%%
\subsection{Global Reduction}
\label{subsec:hpmc-global-reduction}
%%%%%%%%%%%%%%%%%%%%%%%%%%%%%%%%%%%%%%%%%%%%%%%%%%%%%%%%%%%%

A global reduction combines one value from every process.

For example,

\[
\boxed{
s
=
\sum_{j=1}^{p}
s_j.
}
\]

Dot products in distributed iterative solvers require reductions.

Global reductions often require synchronization across the full machine.

%%%%%%%%%%%%%%%%%%%%%%%%%%%%%%%%%%%%%%%%%%%%%%%%%%%%%%%%%%%%
\subsection{Communication Cost Model}
\label{subsec:hpmc-communication-model}
%%%%%%%%%%%%%%%%%%%%%%%%%%%%%%%%%%%%%%%%%%%%%%%%%%%%%%%%%%%%

A simple communication model for sending \(m\) words is

\[
\boxed{
T_{\mathrm{comm}}
=
\alpha+\beta m,
}
\]

where

\[
\alpha
=
\text{latency},
\]

and

\[
\beta
=
\text{time per transmitted word}.
\]

This explains why many tiny messages can be inefficient.

%%%%%%%%%%%%%%%%%%%%%%%%%%%%%%%%%%%%%%%%%%%%%%%%%%%%%%%%%%%%
\subsection{Latency}
\label{subsec:hpmc-latency}
%%%%%%%%%%%%%%%%%%%%%%%%%%%%%%%%%%%%%%%%%%%%%%%%%%%%%%%%%%%%

Latency is the fixed cost of initiating a communication.

Even a very small message incurs this startup cost.

Thus a program issuing millions of tiny messages can be much slower than
one transferring the same total data in larger batches.

%%%%%%%%%%%%%%%%%%%%%%%%%%%%%%%%%%%%%%%%%%%%%%%%%%%%%%%%%%%%
\subsection{Bandwidth}
\label{subsec:hpmc-network-bandwidth}
%%%%%%%%%%%%%%%%%%%%%%%%%%%%%%%%%%%%%%%%%%%%%%%%%%%%%%%%%%%%

Network bandwidth determines the rate at which large messages can be
transferred.

When messages are large, communication cost is dominated increasingly by

\[
\boxed{
\beta m.
}
\]

Both latency and bandwidth must therefore be considered.

%%%%%%%%%%%%%%%%%%%%%%%%%%%%%%%%%%%%%%%%%%%%%%%%%%%%%%%%%%%%
\subsection{Communication Avoidance}
\label{subsec:communication-avoidance}
%%%%%%%%%%%%%%%%%%%%%%%%%%%%%%%%%%%%%%%%%%%%%%%%%%%%%%%%%%%%

A communication-avoiding algorithm reorganizes computation to reduce data
movement or synchronization.

The goal may be to reduce:

\[
\boxed{
\text{number of messages},
}
\]

\[
\boxed{
\text{total communicated data},
}
\]

or

\[
\boxed{
\text{global reductions}.
}
\]

On modern large machines, reducing communication can be more important than
reducing arithmetic operations.

%%%%%%%%%%%%%%%%%%%%%%%%%%%%%%%%%%%%%%%%%%%%%%%%%%%%%%%%%%%%
\subsection{Communication Hiding}
\label{subsec:communication-hiding}
%%%%%%%%%%%%%%%%%%%%%%%%%%%%%%%%%%%%%%%%%%%%%%%%%%%%%%%%%%%%

Communication hiding overlaps communication with useful computation.

Conceptually,

\[
\boxed{
\text{communicate data}
\parallel
\text{compute independent work}.
}
\]

If enough independent work exists, part of the communication time becomes
invisible in total runtime.

%%%%%%%%%%%%%%%%%%%%%%%%%%%%%%%%%%%%%%%%%%%%%%%%%%%%%%%%%%%%
\subsection{Nonblocking Communication}
\label{subsec:nonblocking-communication}
%%%%%%%%%%%%%%%%%%%%%%%%%%%%%%%%%%%%%%%%%%%%%%%%%%%%%%%%%%%%

Nonblocking communication allows a program to initiate a transfer and
continue computing before the transfer completes.

Later, the program checks or waits for completion.

This supports communication overlap but requires careful dependency
management.

%%%%%%%%%%%%%%%%%%%%%%%%%%%%%%%%%%%%%%%%%%%%%%%%%%%%%%%%%%%%
\subsection{Domain Decomposition}
\label{subsec:hpmc-domain-decomposition}
%%%%%%%%%%%%%%%%%%%%%%%%%%%%%%%%%%%%%%%%%%%%%%%%%%%%%%%%%%%%

Suppose

\[
\Omega
=
\Omega_1
\cup
\cdots
\cup
\Omega_p.
\]

Process \(j\) owns subdomain

\[
\Omega_j.
\]

Each process performs local calculations while exchanging data along
subdomain boundaries.

This is a central strategy for scalable PDE simulation.

%%%%%%%%%%%%%%%%%%%%%%%%%%%%%%%%%%%%%%%%%%%%%%%%%%%%%%%%%%%%
\subsection{Halo Exchange}
\label{subsec:hpmc-halo-exchange}
%%%%%%%%%%%%%%%%%%%%%%%%%%%%%%%%%%%%%%%%%%%%%%%%%%%%%%%%%%%%

Suppose a finite-difference stencil requires neighboring values.

Near a subdomain boundary, some neighbors are stored by another process.

Local copies are stored in ghost or halo cells.

At each update,

\[
\boxed{
\text{neighboring processes exchange halo values}.
}
\]

The interior can then be updated locally.

%%%%%%%%%%%%%%%%%%%%%%%%%%%%%%%%%%%%%%%%%%%%%%%%%%%%%%%%%%%%
\subsection{Interior and Boundary Work}
\label{subsec:interior-boundary-work}
%%%%%%%%%%%%%%%%%%%%%%%%%%%%%%%%%%%%%%%%%%%%%%%%%%%%%%%%%%%%

A domain-decomposed PDE step can often be organized as:

\begin{enumerate}
    \item initiate halo communication;
    \item update interior points;
    \item complete communication;
    \item update boundary-adjacent points.
\end{enumerate}

This overlaps local computation with communication.

%%%%%%%%%%%%%%%%%%%%%%%%%%%%%%%%%%%%%%%%%%%%%%%%%%%%%%%%%%%%
\subsection{Surface-to-Volume Scaling}
\label{subsec:hpmc-surface-volume-scaling}
%%%%%%%%%%%%%%%%%%%%%%%%%%%%%%%%%%%%%%%%%%%%%%%%%%%%%%%%%%%%

For a three-dimensional cubic subdomain with side length \(n\),

\[
\boxed{
\text{computation}
\sim
n^3,
}
\]

while interface communication scales approximately as

\[
\boxed{
n^2.
}
\]

As local subdomains become smaller, the communication-to-computation ratio
increases.

This limits strong scaling.

%%%%%%%%%%%%%%%%%%%%%%%%%%%%%%%%%%%%%%%%%%%%%%%%%%%%%%%%%%%%
\subsection{Parallel Speedup}
\label{subsec:hpmc-speedup}
%%%%%%%%%%%%%%%%%%%%%%%%%%%%%%%%%%%%%%%%%%%%%%%%%%%%%%%%%%%%

If serial runtime is

\[
T_1
\]

and runtime using \(p\) processing units is

\[
T_p,
\]

define speedup

\[
\boxed{
S_p
=
\frac{
T_1
}{
T_p
}.
}
\]

Ideal scaling gives

\[
\boxed{
S_p=p.
}
\]

%%%%%%%%%%%%%%%%%%%%%%%%%%%%%%%%%%%%%%%%%%%%%%%%%%%%%%%%%%%%
\subsection{Parallel Efficiency}
\label{subsec:hpmc-parallel-efficiency}
%%%%%%%%%%%%%%%%%%%%%%%%%%%%%%%%%%%%%%%%%%%%%%%%%%%%%%%%%%%%

Parallel efficiency is

\[
\boxed{
E_p
=
\frac{
S_p
}{
p
}.
}
\]

Equivalently,

\[
\boxed{
E_p
=
\frac{
T_1
}{
pT_p
}.
}
\]

Perfect efficiency gives

\[
\boxed{
E_p=1.
}
\]

%%%%%%%%%%%%%%%%%%%%%%%%%%%%%%%%%%%%%%%%%%%%%%%%%%%%%%%%%%%%
\subsection{Worked Example: Parallel Scaling}
\label{subsec:worked-hpmc-scaling}
%%%%%%%%%%%%%%%%%%%%%%%%%%%%%%%%%%%%%%%%%%%%%%%%%%%%%%%%%%%%

\begin{workedexamplebox}{Speedup and Efficiency}

Suppose a computation requires

\[
T_1=800
\]

seconds on one processor.

On \(16\) processors it requires

\[
T_{16}=60
\]

seconds.

Then

\[
S_{16}
=
\frac{800}{60}
\approx
13.33.
\]

The efficiency is

\[
E_{16}
=
\frac{
13.33
}{
16
}
\approx
0.833.
\]

\begin{answerbox}
\[
\boxed{
S_{16}\approx13.33,
\qquad
E_{16}\approx83.3\%.
}
\]
\end{answerbox}

\end{workedexamplebox}

%%%%%%%%%%%%%%%%%%%%%%%%%%%%%%%%%%%%%%%%%%%%%%%%%%%%%%%%%%%%
\subsection{Amdahl's Law}
\label{subsec:hpmc-amdahl}
%%%%%%%%%%%%%%%%%%%%%%%%%%%%%%%%%%%%%%%%%%%%%%%%%%%%%%%%%%%%

Suppose fraction \(s\) of a computation remains serial.

Then idealized speedup is

\[
\boxed{
S_p
=
\frac{
1
}{
s+\frac{1-s}{p}
}.
}
\]

As

\[
p\to\infty,
\]

\[
\boxed{
S_{\max}
=
\frac1s.
}
\]

Thus serial work limits strong scaling.

%%%%%%%%%%%%%%%%%%%%%%%%%%%%%%%%%%%%%%%%%%%%%%%%%%%%%%%%%%%%
\subsection{Worked Example: Serial Fraction}
\label{subsec:worked-serial-fraction}
%%%%%%%%%%%%%%%%%%%%%%%%%%%%%%%%%%%%%%%%%%%%%%%%%%%%%%%%%%%%

\begin{workedexamplebox}{Amdahl Limit}

Suppose

\[
s=0.02.
\]

Then

\[
S_{\max}
=
\frac1{0.02}
=
50.
\]

\begin{answerbox}
\[
\boxed{
S_{\max}=50.
}
\]
\end{answerbox}

Even infinitely many processors cannot overcome the serial fraction in this
idealized model.

\end{workedexamplebox}

%%%%%%%%%%%%%%%%%%%%%%%%%%%%%%%%%%%%%%%%%%%%%%%%%%%%%%%%%%%%
\subsection{Strong Scaling}
\label{subsec:hpmc-strong-scaling}
%%%%%%%%%%%%%%%%%%%%%%%%%%%%%%%%%%%%%%%%%%%%%%%%%%%%%%%%%%%%

Strong scaling keeps total problem size fixed while increasing \(p\).

The goal is to reduce

\[
\boxed{
T_p.
}
\]

Eventually each processor receives too little work relative to
communication and synchronization overhead.

%%%%%%%%%%%%%%%%%%%%%%%%%%%%%%%%%%%%%%%%%%%%%%%%%%%%%%%%%%%%
\subsection{Weak Scaling}
\label{subsec:hpmc-weak-scaling}
%%%%%%%%%%%%%%%%%%%%%%%%%%%%%%%%%%%%%%%%%%%%%%%%%%%%%%%%%%%%

Weak scaling increases total problem size as the number of processors
increases.

If each process performs approximately the same amount of local work, the
ideal weak-scaling runtime remains nearly constant.

Thus weak scaling asks whether additional hardware can support
proportionally larger problems.

%%%%%%%%%%%%%%%%%%%%%%%%%%%%%%%%%%%%%%%%%%%%%%%%%%%%%%%%%%%%
\subsection{Gustafson-Type Scaling}
\label{subsec:gustafson-scaling}
%%%%%%%%%%%%%%%%%%%%%%%%%%%%%%%%%%%%%%%%%%%%%%%%%%%%%%%%%%%%

Amdahl's law assumes a fixed total problem size.

In many scientific applications, adding processors allows the problem to
grow.

Gustafson-type reasoning emphasizes that the parallel fraction can increase
with problem size.

Thus large machines may be useful even when strong-scaling speedup is
limited.

%%%%%%%%%%%%%%%%%%%%%%%%%%%%%%%%%%%%%%%%%%%%%%%%%%%%%%%%%%%%
\subsection{Load Imbalance}
\label{subsec:hpmc-load-imbalance}
%%%%%%%%%%%%%%%%%%%%%%%%%%%%%%%%%%%%%%%%%%%%%%%%%%%%%%%%%%%%

Suppose worker \(j\) requires time

\[
T_j.
\]

If all workers synchronize at the end, total time is approximately

\[
\boxed{
\max_j T_j.
}
\]

Thus one slow worker can delay every other worker.

Balanced partitioning seeks

\[
\boxed{
T_1
\approx
T_2
\approx
\cdots
\approx
T_p.
}
\]

%%%%%%%%%%%%%%%%%%%%%%%%%%%%%%%%%%%%%%%%%%%%%%%%%%%%%%%%%%%%
\subsection{Static and Dynamic Load Balancing}
\label{subsec:static-dynamic-load-balancing}
%%%%%%%%%%%%%%%%%%%%%%%%%%%%%%%%%%%%%%%%%%%%%%%%%%%%%%%%%%%%

Static load balancing assigns tasks before execution.

Dynamic load balancing redistributes work while computation proceeds.

Static assignment has lower overhead.

Dynamic assignment is useful when task costs are difficult to predict.

%%%%%%%%%%%%%%%%%%%%%%%%%%%%%%%%%%%%%%%%%%%%%%%%%%%%%%%%%%%%
\subsection{Accelerator Computing}
\label{subsec:accelerator-computing}
%%%%%%%%%%%%%%%%%%%%%%%%%%%%%%%%%%%%%%%%%%%%%%%%%%%%%%%%%%%%

Accelerators are specialized processors designed for high computational
throughput.

The most common examples are GPUs.

They are especially effective for:

\[
\boxed{
\text{regular parallel arithmetic},
}
\]

\[
\boxed{
\text{dense linear algebra},
}
\]

and

\[
\boxed{
\text{large independent batches}.
}
\]

%%%%%%%%%%%%%%%%%%%%%%%%%%%%%%%%%%%%%%%%%%%%%%%%%%%%%%%%%%%%
\subsection{GPU Parallelism}
\label{subsec:gpu-parallelism}
%%%%%%%%%%%%%%%%%%%%%%%%%%%%%%%%%%%%%%%%%%%%%%%%%%%%%%%%%%%%

GPUs execute many lightweight computational threads.

Performance is best when many threads perform similar operations on large
amounts of data.

Examples include:

\[
\boxed{
y_i=ax_i+b_i,
}
\]

and grid-based simulation updates.

%%%%%%%%%%%%%%%%%%%%%%%%%%%%%%%%%%%%%%%%%%%%%%%%%%%%%%%%%%%%
\subsection{Host and Device}
\label{subsec:host-device}
%%%%%%%%%%%%%%%%%%%%%%%%%%%%%%%%%%%%%%%%%%%%%%%%%%%%%%%%%%%%

In a heterogeneous system, one often distinguishes:

\[
\boxed{
\text{host}
=
\text{CPU side},
}
\]

and

\[
\boxed{
\text{device}
=
\text{accelerator side}.
}
\]

Data may need to move between separate memory spaces.

Such transfers can be costly.

%%%%%%%%%%%%%%%%%%%%%%%%%%%%%%%%%%%%%%%%%%%%%%%%%%%%%%%%%%%%
\subsection{Accelerator Transfer Cost}
\label{subsec:accelerator-transfer-cost}
%%%%%%%%%%%%%%%%%%%%%%%%%%%%%%%%%%%%%%%%%%%%%%%%%%%%%%%%%%%%

Suppose computation takes

\[
T_{\mathrm{kernel}}
\]

but transferring input and output requires

\[
T_{\mathrm{transfer}}.
\]

Total accelerator time is approximately

\[
\boxed{
T
=
T_{\mathrm{transfer}}
+
T_{\mathrm{kernel}}.
}
\]

If

\[
T_{\mathrm{transfer}}
\gg
T_{\mathrm{kernel}},
\]

the accelerator may provide little overall benefit.

%%%%%%%%%%%%%%%%%%%%%%%%%%%%%%%%%%%%%%%%%%%%%%%%%%%%%%%%%%%%
\subsection{Kernel Launch Overhead}
\label{subsec:kernel-launch-overhead}
%%%%%%%%%%%%%%%%%%%%%%%%%%%%%%%%%%%%%%%%%%%%%%%%%%%%%%%%%%%%

Starting an accelerator kernel incurs a fixed overhead.

Therefore thousands of tiny kernels may be inefficient.

Combining operations into larger kernels can reduce launch overhead and
intermediate memory traffic.

%%%%%%%%%%%%%%%%%%%%%%%%%%%%%%%%%%%%%%%%%%%%%%%%%%%%%%%%%%%%
\subsection{Kernel Fusion}
\label{subsec:kernel-fusion}
%%%%%%%%%%%%%%%%%%%%%%%%%%%%%%%%%%%%%%%%%%%%%%%%%%%%%%%%%%%%

Suppose an algorithm computes

\[
y_i=f(x_i)
\]

and then

\[
z_i=g(y_i).
\]

Instead of storing all \(y_i\), a fused kernel can compute

\[
\boxed{
z_i
=
g(f(x_i))
}
\]

directly.

This reduces:

\[
\boxed{
\text{memory traffic},
}
\]

and

\[
\boxed{
\text{kernel launches}.
}
\]

%%%%%%%%%%%%%%%%%%%%%%%%%%%%%%%%%%%%%%%%%%%%%%%%%%%%%%%%%%%%
\subsection{Batch Computation}
\label{subsec:batch-computation}
%%%%%%%%%%%%%%%%%%%%%%%%%%%%%%%%%%%%%%%%%%%%%%%%%%%%%%%%%%%%

Many small operations can be grouped into a larger batch.

Examples include:

\[
\boxed{
\text{many small matrix factorizations},
}
\]

\[
\boxed{
\text{many Monte Carlo trajectories},
}
\]

and

\[
\boxed{
\text{many independent model evaluations}.
}
\]

Batching increases available parallelism and amortizes overhead.

%%%%%%%%%%%%%%%%%%%%%%%%%%%%%%%%%%%%%%%%%%%%%%%%%%%%%%%%%%%%
\subsection{Mixed Precision}
\label{subsec:hpmc-mixed-precision}
%%%%%%%%%%%%%%%%%%%%%%%%%%%%%%%%%%%%%%%%%%%%%%%%%%%%%%%%%%%%

Different parts of a computation may use different numerical precisions.

For example,

\[
\boxed{
\text{low precision}
}
\]

may be used for an approximate solve, followed by

\[
\boxed{
\text{high-precision correction}.
}
\]

This can reduce memory traffic and increase arithmetic throughput.

%%%%%%%%%%%%%%%%%%%%%%%%%%%%%%%%%%%%%%%%%%%%%%%%%%%%%%%%%%%%
\subsection{Iterative Refinement}
\label{subsec:hpmc-iterative-refinement}
%%%%%%%%%%%%%%%%%%%%%%%%%%%%%%%%%%%%%%%%%%%%%%%%%%%%%%%%%%%%

Suppose an approximate solution \(x_k\) of

\[
Ax=b
\]

is available.

Compute high-accuracy residual

\[
\boxed{
r_k
=
b-Ax_k.
}
\]

Solve approximately

\[
\boxed{
A\delta x_k
=
r_k
}
\]

and update

\[
\boxed{
x_{k+1}
=
x_k+\delta x_k.
}
\]

The inner solve may sometimes use lower precision than the residual
calculation.

%%%%%%%%%%%%%%%%%%%%%%%%%%%%%%%%%%%%%%%%%%%%%%%%%%%%%%%%%%%%
\subsection{Mixed-Precision Tradeoff}
\label{subsec:mixed-precision-tradeoff}
%%%%%%%%%%%%%%%%%%%%%%%%%%%%%%%%%%%%%%%%%%%%%%%%%%%%%%%%%%%%

Lower precision can improve:

\[
\boxed{
\text{throughput},
}
\]

and

\[
\boxed{
\text{memory efficiency}.
}
\]

But it may worsen:

\[
\boxed{
\text{rounding error},
}
\]

and

\[
\boxed{
\text{solver convergence}.
}
\]

Conditioning determines whether lower precision is acceptable.

%%%%%%%%%%%%%%%%%%%%%%%%%%%%%%%%%%%%%%%%%%%%%%%%%%%%%%%%%%%%
\subsection{Distributed Dense Matrices}
\label{subsec:distributed-dense-matrices}
%%%%%%%%%%%%%%%%%%%%%%%%%%%%%%%%%%%%%%%%%%%%%%%%%%%%%%%%%%%%

A large dense matrix may be partitioned into blocks across processes.

For example,

\[
\boxed{
A
=
\begin{pmatrix}
A_{11}&A_{12}\\
A_{21}&A_{22}
\end{pmatrix}.
}
\]

Each process owns one or more blocks.

Parallel factorization requires both local matrix operations and
communication of pivot or panel information.

%%%%%%%%%%%%%%%%%%%%%%%%%%%%%%%%%%%%%%%%%%%%%%%%%%%%%%%%%%%%
\subsection{One-Dimensional Matrix Distribution}
\label{subsec:one-dimensional-matrix-distribution}
%%%%%%%%%%%%%%%%%%%%%%%%%%%%%%%%%%%%%%%%%%%%%%%%%%%%%%%%%%%%

A simple parallel layout assigns blocks of rows to processes.

Process \(j\) owns

\[
\boxed{
A_{I_j,:}.
}
\]

This is easy to implement but can create communication bottlenecks for some
matrix operations.

%%%%%%%%%%%%%%%%%%%%%%%%%%%%%%%%%%%%%%%%%%%%%%%%%%%%%%%%%%%%
\subsection{Two-Dimensional Matrix Distribution}
\label{subsec:two-dimensional-matrix-distribution}
%%%%%%%%%%%%%%%%%%%%%%%%%%%%%%%%%%%%%%%%%%%%%%%%%%%%%%%%%%%%

A more scalable dense layout distributes matrix blocks across a
two-dimensional processor grid.

This reduces communication per processor for many matrix factorizations and
matrix multiplications.

The matrix distribution becomes part of the algorithm.

%%%%%%%%%%%%%%%%%%%%%%%%%%%%%%%%%%%%%%%%%%%%%%%%%%%%%%%%%%%%
\subsection{Parallel Matrix Multiplication}
\label{subsec:parallel-matrix-multiplication}
%%%%%%%%%%%%%%%%%%%%%%%%%%%%%%%%%%%%%%%%%%%%%%%%%%%%%%%%%%%%

For

\[
C=AB,
\]

processes compute local block products and communicate required blocks.

A scalable implementation tries to balance:

\[
\boxed{
\text{local arithmetic}
}
\]

and

\[
\boxed{
\text{data exchange}.
}
\]

Because matrix multiplication has high arithmetic intensity, it scales
well on many architectures.

%%%%%%%%%%%%%%%%%%%%%%%%%%%%%%%%%%%%%%%%%%%%%%%%%%%%%%%%%%%%
\subsection{Parallel Factorizations}
\label{subsec:parallel-factorizations}
%%%%%%%%%%%%%%%%%%%%%%%%%%%%%%%%%%%%%%%%%%%%%%%%%%%%%%%%%%%%

Parallel versions of:

\[
\boxed{
\text{LU},
}
\]

\[
\boxed{
\text{QR},
}
\]

and

\[
\boxed{
\text{Cholesky}
}
\]

divide the matrix into blocks.

Large matrix-matrix updates are performed efficiently, while smaller panel
or synchronization steps may limit scalability.

%%%%%%%%%%%%%%%%%%%%%%%%%%%%%%%%%%%%%%%%%%%%%%%%%%%%%%%%%%%%
\subsection{Communication-Avoiding QR Preview}
\label{subsec:communication-avoiding-qr}
%%%%%%%%%%%%%%%%%%%%%%%%%%%%%%%%%%%%%%%%%%%%%%%%%%%%%%%%%%%%

Classical QR factorization can require repeated communication.

Communication-avoiding QR reorganizes the factorization to reduce data
movement while preserving mathematical stability.

This is an example where the numerical algorithm itself is redesigned for
modern hardware.

%%%%%%%%%%%%%%%%%%%%%%%%%%%%%%%%%%%%%%%%%%%%%%%%%%%%%%%%%%%%
\subsection{Parallel Krylov Methods}
\label{subsec:parallel-krylov-methods}
%%%%%%%%%%%%%%%%%%%%%%%%%%%%%%%%%%%%%%%%%%%%%%%%%%%%%%%%%%%%

Krylov methods rely on operations such as:

\[
\boxed{
Ax,
}
\]

\[
\boxed{
x^Ty,
}
\]

and

\[
\boxed{
\|x\|.
}
\]

Sparse matrix-vector products often require local neighbor communication.

Dot products require global reductions.

Thus communication can dominate at large processor counts.

%%%%%%%%%%%%%%%%%%%%%%%%%%%%%%%%%%%%%%%%%%%%%%%%%%%%%%%%%%%%
\subsection{Global Reduction Bottleneck}
\label{subsec:krylov-global-reduction}
%%%%%%%%%%%%%%%%%%%%%%%%%%%%%%%%%%%%%%%%%%%%%%%%%%%%%%%%%%%%

A dot product

\[
x^Ty
=
\sum_i x_i y_i
\]

is computed locally first.

Then all processes combine their partial sums.

This global reduction requires collective synchronization.

At extreme scale, the cost of a few scalar reductions can exceed substantial
local arithmetic.

%%%%%%%%%%%%%%%%%%%%%%%%%%%%%%%%%%%%%%%%%%%%%%%%%%%%%%%%%%%%
\subsection{Pipelined Krylov Methods Preview}
\label{subsec:pipelined-krylov}
%%%%%%%%%%%%%%%%%%%%%%%%%%%%%%%%%%%%%%%%%%%%%%%%%%%%%%%%%%%%

Pipelined Krylov methods reorganize iterations so communication associated
with reductions can overlap with matrix-vector computations.

The mathematical recurrence becomes more complicated.

The benefit is potentially reduced synchronization cost.

%%%%%%%%%%%%%%%%%%%%%%%%%%%%%%%%%%%%%%%%%%%%%%%%%%%%%%%%%%%%
\subsection{Preconditioning at Scale}
\label{subsec:parallel-preconditioning}
%%%%%%%%%%%%%%%%%%%%%%%%%%%%%%%%%%%%%%%%%%%%%%%%%%%%%%%%%%%%

A preconditioner must improve convergence while also remaining parallel.

A theoretically excellent serial preconditioner may scale poorly.

Parallel-friendly choices include:

\[
\boxed{
\text{block methods},
}
\]

\[
\boxed{
\text{domain decomposition},
}
\]

and

\[
\boxed{
\text{multigrid}.
}
\]

%%%%%%%%%%%%%%%%%%%%%%%%%%%%%%%%%%%%%%%%%%%%%%%%%%%%%%%%%%%%
\subsection{Parallel Multigrid}
\label{subsec:parallel-multigrid}
%%%%%%%%%%%%%%%%%%%%%%%%%%%%%%%%%%%%%%%%%%%%%%%%%%%%%%%%%%%%

Multigrid combines work on several mesh resolutions.

Fine grids provide abundant parallel work.

Coarser grids contain fewer unknowns and eventually provide insufficient
parallelism.

Thus coarse-grid solves can become scalability bottlenecks.

%%%%%%%%%%%%%%%%%%%%%%%%%%%%%%%%%%%%%%%%%%%%%%%%%%%%%%%%%%%%
\subsection{Coarse-Grid Bottleneck}
\label{subsec:multigrid-coarse-grid-bottleneck}
%%%%%%%%%%%%%%%%%%%%%%%%%%%%%%%%%%%%%%%%%%%%%%%%%%%%%%%%%%%%

Suppose the finest level has millions of unknowns distributed across many
processes.

After repeated coarsening, the coarse problem may contain fewer unknowns
than processors.

Then many processors have little or no useful work.

Special coarse-grid strategies are required at large scale.

%%%%%%%%%%%%%%%%%%%%%%%%%%%%%%%%%%%%%%%%%%%%%%%%%%%%%%%%%%%%
\subsection{Parallel Fast Fourier Transform}
\label{subsec:parallel-fft}
%%%%%%%%%%%%%%%%%%%%%%%%%%%%%%%%%%%%%%%%%%%%%%%%%%%%%%%%%%%%

The FFT requires approximately

\[
\boxed{
O(N\log N)
}
\]

arithmetic operations.

In distributed systems, data must also be rearranged among processes.

For multidimensional FFTs, these redistributions can require expensive
global communication.

%%%%%%%%%%%%%%%%%%%%%%%%%%%%%%%%%%%%%%%%%%%%%%%%%%%%%%%%%%%%
\subsection{FFT Communication}
\label{subsec:fft-communication}
%%%%%%%%%%%%%%%%%%%%%%%%%%%%%%%%%%%%%%%%%%%%%%%%%%%%%%%%%%%%

A three-dimensional distributed FFT may partition data into slabs or
pencils.

Transforms are performed locally along one direction.

Then data are redistributed so transforms can be performed along another
direction.

These global transpositions often dominate runtime at large scale.

%%%%%%%%%%%%%%%%%%%%%%%%%%%%%%%%%%%%%%%%%%%%%%%%%%%%%%%%%%%%
\subsection{High-Performance PDE Solvers}
\label{subsec:high-performance-pde-solvers}
%%%%%%%%%%%%%%%%%%%%%%%%%%%%%%%%%%%%%%%%%%%%%%%%%%%%%%%%%%%%

A scalable PDE solver combines:

\[
\boxed{
\text{domain decomposition},
}
\]

\[
\boxed{
\text{local numerical kernels},
}
\]

\[
\boxed{
\text{halo communication},
}
\]

and

\[
\boxed{
\text{scalable linear solvers}.
}
\]

The numerical discretization and parallel decomposition should be designed
together.

%%%%%%%%%%%%%%%%%%%%%%%%%%%%%%%%%%%%%%%%%%%%%%%%%%%%%%%%%%%%
\subsection{Structured-Grid PDEs}
\label{subsec:hpmc-structured-grid-pdes}
%%%%%%%%%%%%%%%%%%%%%%%%%%%%%%%%%%%%%%%%%%%%%%%%%%%%%%%%%%%%

Structured grids have regular communication patterns.

Neighboring processes exchange values across predictable boundaries.

These methods often provide:

\[
\boxed{
\text{simple memory access},
}
\]

\[
\boxed{
\text{efficient vectorization},
}
\]

and

\[
\boxed{
\text{regular parallel decomposition}.
}
\]

%%%%%%%%%%%%%%%%%%%%%%%%%%%%%%%%%%%%%%%%%%%%%%%%%%%%%%%%%%%%
\subsection{Unstructured-Grid PDEs}
\label{subsec:hpmc-unstructured-grid-pdes}
%%%%%%%%%%%%%%%%%%%%%%%%%%%%%%%%%%%%%%%%%%%%%%%%%%%%%%%%%%%%

Unstructured meshes allow complex geometry.

However, they produce irregular:

\[
\boxed{
\text{memory access},
}
\]

\[
\boxed{
\text{communication patterns},
}
\]

and

\[
\boxed{
\text{load balancing}.
}
\]

Graph partitioning becomes important in assigning mesh elements to
processes.

%%%%%%%%%%%%%%%%%%%%%%%%%%%%%%%%%%%%%%%%%%%%%%%%%%%%%%%%%%%%
\subsection{Graph Partitioning}
\label{subsec:hpmc-graph-partitioning}
%%%%%%%%%%%%%%%%%%%%%%%%%%%%%%%%%%%%%%%%%%%%%%%%%%%%%%%%%%%%

Represent a computational mesh as graph

\[
G=(V,E).
\]

Assign vertices to \(p\) partitions while seeking:

\[
\boxed{
\text{balanced vertex counts}
}
\]

and

\[
\boxed{
\text{few edges crossing partitions}.
}
\]

This balances computation while reducing communication.

%%%%%%%%%%%%%%%%%%%%%%%%%%%%%%%%%%%%%%%%%%%%%%%%%%%%%%%%%%%%
\subsection{Adaptive Meshes at Scale}
\label{subsec:hpmc-adaptive-meshes}
%%%%%%%%%%%%%%%%%%%%%%%%%%%%%%%%%%%%%%%%%%%%%%%%%%%%%%%%%%%%

Adaptive refinement creates nonuniform workloads.

As the mesh changes, a previously balanced partition may become highly
imbalanced.

Parallel adaptive methods therefore require periodic repartitioning or
dynamic load balancing.

%%%%%%%%%%%%%%%%%%%%%%%%%%%%%%%%%%%%%%%%%%%%%%%%%%%%%%%%%%%%
\subsection{Parallel Monte Carlo}
\label{subsec:hpmc-parallel-monte-carlo}
%%%%%%%%%%%%%%%%%%%%%%%%%%%%%%%%%%%%%%%%%%%%%%%%%%%%%%%%%%%%

Monte Carlo simulation is naturally parallel.

Suppose

\[
\widehat{\mu}
=
\frac1N
\sum_{i=1}^{N}
g(X_i).
\]

Each worker generates a subset of samples independently.

At the end, partial sums are reduced.

Communication is minimal relative to computation.

%%%%%%%%%%%%%%%%%%%%%%%%%%%%%%%%%%%%%%%%%%%%%%%%%%%%%%%%%%%%
\subsection{Random Number Streams}
\label{subsec:hpmc-random-streams}
%%%%%%%%%%%%%%%%%%%%%%%%%%%%%%%%%%%%%%%%%%%%%%%%%%%%%%%%%%%%

Parallel Monte Carlo requires carefully separated pseudorandom streams.

Workers should not accidentally generate overlapping or strongly
correlated sequences.

Stream management should support both:

\[
\boxed{
\text{statistical quality}
}
\]

and

\[
\boxed{
\text{reproducibility}.
}
\]

%%%%%%%%%%%%%%%%%%%%%%%%%%%%%%%%%%%%%%%%%%%%%%%%%%%%%%%%%%%%
\subsection{Parallel Ensembles}
\label{subsec:hpmc-parallel-ensembles}
%%%%%%%%%%%%%%%%%%%%%%%%%%%%%%%%%%%%%%%%%%%%%%%%%%%%%%%%%%%%

Suppose a simulation must be run for parameter values

\[
\theta_1,\ldots,\theta_M.
\]

Each parameter case can often be assigned independently.

Thus ensemble simulation provides coarse-grained parallelism even when one
individual simulation does not scale efficiently.

%%%%%%%%%%%%%%%%%%%%%%%%%%%%%%%%%%%%%%%%%%%%%%%%%%%%%%%%%%%%
\subsection{Nested Parallelism}
\label{subsec:nested-parallelism}
%%%%%%%%%%%%%%%%%%%%%%%%%%%%%%%%%%%%%%%%%%%%%%%%%%%%%%%%%%%%

A computational study may have parallelism at several levels.

For example:

\[
\boxed{
\text{parameter cases}
}
\]

may be distributed across groups of processors, while each individual case
uses parallel PDE computation internally.

This is nested parallelism.

%%%%%%%%%%%%%%%%%%%%%%%%%%%%%%%%%%%%%%%%%%%%%%%%%%%%%%%%%%%%
\subsection{Parallel Optimization}
\label{subsec:hpmc-parallel-optimization}
%%%%%%%%%%%%%%%%%%%%%%%%%%%%%%%%%%%%%%%%%%%%%%%%%%%%%%%%%%%%

Optimization may exploit parallelism through:

\[
\boxed{
\text{parallel objective evaluations},
}
\]

\[
\boxed{
\text{parallel gradient computations},
}
\]

\[
\boxed{
\text{distributed linear algebra},
}
\]

and

\[
\boxed{
\text{multiple starting points}.
}
\]

The useful strategy depends on the optimization algorithm.

%%%%%%%%%%%%%%%%%%%%%%%%%%%%%%%%%%%%%%%%%%%%%%%%%%%%%%%%%%%%
\subsection{Parallel Gradient Computation}
\label{subsec:parallel-gradient-computation}
%%%%%%%%%%%%%%%%%%%%%%%%%%%%%%%%%%%%%%%%%%%%%%%%%%%%%%%%%%%%

If

\[
F(x)
=
\sum_{i=1}^{N}
f_i(x),
\]

then

\[
\boxed{
\nabla F(x)
=
\sum_{i=1}^{N}
\nabla f_i(x).
}
\]

Each worker can compute a partial gradient.

A global reduction forms the total gradient.

This is common in large data-fitting problems.

%%%%%%%%%%%%%%%%%%%%%%%%%%%%%%%%%%%%%%%%%%%%%%%%%%%%%%%%%%%%
\subsection{Distributed Hessian-Vector Products}
\label{subsec:distributed-hessian-vector}
%%%%%%%%%%%%%%%%%%%%%%%%%%%%%%%%%%%%%%%%%%%%%%%%%%%%%%%%%%%%

For

\[
F(x)
=
\sum_i f_i(x),
\]

the Hessian-vector product is

\[
\boxed{
\nabla^2F(x)v
=
\sum_i
\nabla^2f_i(x)v.
}
\]

Local contributions can be computed in parallel and summed.

This supports large-scale Newton--Krylov optimization without forming the
full Hessian.

%%%%%%%%%%%%%%%%%%%%%%%%%%%%%%%%%%%%%%%%%%%%%%%%%%%%%%%%%%%%
\subsection{Task-Based Parallelism}
\label{subsec:task-based-parallelism}
%%%%%%%%%%%%%%%%%%%%%%%%%%%%%%%%%%%%%%%%%%%%%%%%%%%%%%%%%%%%

Instead of explicitly assigning work to processors, a program can express a
collection of tasks and dependencies.

A runtime system schedules ready tasks dynamically.

This can improve:

\[
\boxed{
\text{load balance},
}
\]

and

\[
\boxed{
\text{overlap of computation and communication}.
}
\]

%%%%%%%%%%%%%%%%%%%%%%%%%%%%%%%%%%%%%%%%%%%%%%%%%%%%%%%%%%%%
\subsection{Dependency Graph}
\label{subsec:hpmc-dependency-graph}
%%%%%%%%%%%%%%%%%%%%%%%%%%%%%%%%%%%%%%%%%%%%%%%%%%%%%%%%%%%%

A task graph is a directed acyclic graph in which nodes represent
computations and edges represent dependencies.

If two tasks have no dependency path between them, they may execute
simultaneously.

This representation exposes algorithmic parallelism.

%%%%%%%%%%%%%%%%%%%%%%%%%%%%%%%%%%%%%%%%%%%%%%%%%%%%%%%%%%%%
\subsection{Critical Path}
\label{subsec:hpmc-critical-path}
%%%%%%%%%%%%%%%%%%%%%%%%%%%%%%%%%%%%%%%%%%%%%%%%%%%%%%%%%%%%

The critical path is the longest dependency chain in a task graph.

Even with infinitely many processors, runtime cannot be shorter than the
critical-path time.

Thus theoretical parallelism is limited by both:

\[
\boxed{
\text{total work}
}
\]

and

\[
\boxed{
\text{dependency depth}.
}
\]

%%%%%%%%%%%%%%%%%%%%%%%%%%%%%%%%%%%%%%%%%%%%%%%%%%%%%%%%%%%%
\subsection{Work and Span}
\label{subsec:work-span}
%%%%%%%%%%%%%%%%%%%%%%%%%%%%%%%%%%%%%%%%%%%%%%%%%%%%%%%%%%%%

Let

\[
T_1
\]

denote total serial work and

\[
T_\infty
\]

denote ideal critical-path time with unlimited processors.

Then the maximum theoretical parallelism is approximately

\[
\boxed{
\frac{
T_1
}{
T_\infty
}.
}
\]

This provides an algorithmic view of parallel scalability.

%%%%%%%%%%%%%%%%%%%%%%%%%%%%%%%%%%%%%%%%%%%%%%%%%%%%%%%%%%%%
\subsection{Asynchronous Iteration}
\label{subsec:hpmc-asynchronous-iteration}
%%%%%%%%%%%%%%%%%%%%%%%%%%%%%%%%%%%%%%%%%%%%%%%%%%%%%%%%%%%%

Synchronous iterative methods require all workers to advance together.

Asynchronous methods allow updates using partially delayed information.

This can reduce waiting and improve hardware utilization.

However, convergence analysis becomes more complicated.

%%%%%%%%%%%%%%%%%%%%%%%%%%%%%%%%%%%%%%%%%%%%%%%%%%%%%%%%%%%%
\subsection{Fault Tolerance}
\label{subsec:hpmc-fault-tolerance}
%%%%%%%%%%%%%%%%%%%%%%%%%%%%%%%%%%%%%%%%%%%%%%%%%%%%%%%%%%%%

A computation running across thousands of processors for many hours has a
greater chance of encountering hardware or software failure.

Fault tolerance allows computation to continue or restart efficiently.

Common approaches include:

\[
\boxed{
\text{checkpoint/restart},
}
\]

\[
\boxed{
\text{replication},
}
\]

and

\[
\boxed{
\text{algorithmic recovery}.
}
\]

%%%%%%%%%%%%%%%%%%%%%%%%%%%%%%%%%%%%%%%%%%%%%%%%%%%%%%%%%%%%
\subsection{Checkpointing at Scale}
\label{subsec:hpmc-checkpointing}
%%%%%%%%%%%%%%%%%%%%%%%%%%%%%%%%%%%%%%%%%%%%%%%%%%%%%%%%%%%%

Suppose a distributed simulation stores state

\[
x_j
\]

on process \(j\).

A global checkpoint saves enough distributed state to reconstruct

\[
\boxed{
x
=
(x_1,\ldots,x_p).
}
\]

Writing checkpoints too frequently can overwhelm storage systems.

Writing them too rarely increases recomputation after failure.

%%%%%%%%%%%%%%%%%%%%%%%%%%%%%%%%%%%%%%%%%%%%%%%%%%%%%%%%%%%%
\subsection{Incremental Checkpointing}
\label{subsec:incremental-checkpointing}
%%%%%%%%%%%%%%%%%%%%%%%%%%%%%%%%%%%%%%%%%%%%%%%%%%%%%%%%%%%%

Instead of writing the entire simulation state each time, an incremental
checkpoint records only changed or necessary portions.

This can reduce storage traffic.

The restart procedure becomes more complex because several checkpoint
pieces may need to be combined.

%%%%%%%%%%%%%%%%%%%%%%%%%%%%%%%%%%%%%%%%%%%%%%%%%%%%%%%%%%%%
\subsection{Resilience}
\label{subsec:hpmc-resilience}
%%%%%%%%%%%%%%%%%%%%%%%%%%%%%%%%%%%%%%%%%%%%%%%%%%%%%%%%%%%%

Resilience is the ability of a computational method to tolerate and recover
from failures or perturbations.

At extreme scale, resilience may need to be part of algorithm design rather
than only an operating-system feature.

%%%%%%%%%%%%%%%%%%%%%%%%%%%%%%%%%%%%%%%%%%%%%%%%%%%%%%%%%%%%
\subsection{Floating-Point Reproducibility}
\label{subsec:hpmc-floating-point-reproducibility}
%%%%%%%%%%%%%%%%%%%%%%%%%%%%%%%%%%%%%%%%%%%%%%%%%%%%%%%%%%%%

Floating-point addition is not associative:

\[
\boxed{
(a+b)+c
\neq
a+(b+c)
}
\]

in general.

Parallel reductions may use different summation trees.

Therefore the final low-order digits can vary with:

\[
\boxed{
\text{processor count},
}
\]

\[
\boxed{
\text{execution order},
}
\]

or

\[
\boxed{
\text{hardware}.
}
\]

%%%%%%%%%%%%%%%%%%%%%%%%%%%%%%%%%%%%%%%%%%%%%%%%%%%%%%%%%%%%
\subsection{Bitwise Versus Numerical Reproducibility}
\label{subsec:hpmc-bitwise-numerical-reproducibility}
%%%%%%%%%%%%%%%%%%%%%%%%%%%%%%%%%%%%%%%%%%%%%%%%%%%%%%%%%%%%

Bitwise reproducibility requires exactly identical binary results.

Numerical reproducibility requires results to agree within meaningful
mathematical tolerances.

For many scientific applications,

\[
\boxed{
\text{numerical reproducibility}
}
\]

is more important than exact bitwise identity.

%%%%%%%%%%%%%%%%%%%%%%%%%%%%%%%%%%%%%%%%%%%%%%%%%%%%%%%%%%%%
\subsection{Reproducible Reductions}
\label{subsec:reproducible-reductions}
%%%%%%%%%%%%%%%%%%%%%%%%%%%%%%%%%%%%%%%%%%%%%%%%%%%%%%%%%%%%

More reproducible summation may use:

\[
\boxed{
\text{fixed reduction trees},
}
\]

\[
\boxed{
\text{pairwise summation},
}
\]

or

\[
\boxed{
\text{compensated accumulation}.
}
\]

These methods may introduce additional computational cost.

Thus reproducibility itself can become a performance tradeoff.

%%%%%%%%%%%%%%%%%%%%%%%%%%%%%%%%%%%%%%%%%%%%%%%%%%%%%%%%%%%%
\subsection{Profiling at Scale}
\label{subsec:hpmc-profiling}
%%%%%%%%%%%%%%%%%%%%%%%%%%%%%%%%%%%%%%%%%%%%%%%%%%%%%%%%%%%%

Large parallel applications should be profiled to identify:

\[
\boxed{
\text{computational hotspots},
}
\]

\[
\boxed{
\text{communication bottlenecks},
}
\]

\[
\boxed{
\text{load imbalance},
}
\]

and

\[
\boxed{
\text{memory limitations}.
}
\]

Optimization should target the dominant measured costs.

%%%%%%%%%%%%%%%%%%%%%%%%%%%%%%%%%%%%%%%%%%%%%%%%%%%%%%%%%%%%
\subsection{Parallel Timeline Analysis}
\label{subsec:parallel-timeline-analysis}
%%%%%%%%%%%%%%%%%%%%%%%%%%%%%%%%%%%%%%%%%%%%%%%%%%%%%%%%%%%%

A parallel timeline shows when each process is:

\[
\boxed{
\text{computing},
}
\]

\[
\boxed{
\text{communicating},
}
\]

or

\[
\boxed{
\text{waiting}.
}
\]

Large idle regions often reveal synchronization or load-balance problems.

%%%%%%%%%%%%%%%%%%%%%%%%%%%%%%%%%%%%%%%%%%%%%%%%%%%%%%%%%%%%
\subsection{Scaling Experiments}
\label{subsec:hpmc-scaling-experiments}
%%%%%%%%%%%%%%%%%%%%%%%%%%%%%%%%%%%%%%%%%%%%%%%%%%%%%%%%%%%%

A strong-scaling experiment records

\[
T_1,T_2,T_4,T_8,\ldots
\]

for one fixed problem.

Then compute:

\[
\boxed{
S_p
=
\frac{T_1}{T_p},
}
\]

and

\[
\boxed{
E_p
=
\frac{S_p}{p}.
}
\]

A weak-scaling experiment increases problem size with \(p\).

%%%%%%%%%%%%%%%%%%%%%%%%%%%%%%%%%%%%%%%%%%%%%%%%%%%%%%%%%%%%
\subsection{Efficiency Loss Decomposition}
\label{subsec:efficiency-loss-decomposition}
%%%%%%%%%%%%%%%%%%%%%%%%%%%%%%%%%%%%%%%%%%%%%%%%%%%%%%%%%%%%

Parallel inefficiency can be attributed conceptually to:

\[
\boxed{
\text{serial work},
}
\]

\[
\boxed{
\text{load imbalance},
}
\]

\[
\boxed{
\text{communication},
}
\]

\[
\boxed{
\text{synchronization},
}
\]

and

\[
\boxed{
\text{resource contention}.
}
\]

Understanding which contribution dominates guides optimization.

%%%%%%%%%%%%%%%%%%%%%%%%%%%%%%%%%%%%%%%%%%%%%%%%%%%%%%%%%%%%
\subsection{Energy Efficiency}
\label{subsec:hpmc-energy-efficiency}
%%%%%%%%%%%%%%%%%%%%%%%%%%%%%%%%%%%%%%%%%%%%%%%%%%%%%%%%%%%%

High-performance computation consumes electrical energy.

An algorithm may be evaluated by quantities such as

\[
\boxed{
\frac{
\text{useful computation}
}{
\text{energy consumed}
}.
}
\]

Reducing data movement can improve both runtime and energy efficiency.

Thus performance and energy objectives often align.

%%%%%%%%%%%%%%%%%%%%%%%%%%%%%%%%%%%%%%%%%%%%%%%%%%%%%%%%%%%%
\subsection{Time to Solution}
\label{subsec:time-to-solution}
%%%%%%%%%%%%%%%%%%%%%%%%%%%%%%%%%%%%%%%%%%%%%%%%%%%%%%%%%%%%

The most relevant performance metric is often total time required to obtain
a scientifically acceptable result.

This includes:

\[
\boxed{
\text{setup},
}
\]

\[
\boxed{
\text{computation},
}
\]

\[
\boxed{
\text{communication},
}
\]

\[
\boxed{
\text{I/O},
}
\]

and

\[
\boxed{
\text{convergence}.
}
\]

A method achieving higher FLOP/s may still have worse time to solution.

%%%%%%%%%%%%%%%%%%%%%%%%%%%%%%%%%%%%%%%%%%%%%%%%%%%%%%%%%%%%
\subsection{Accuracy to Solution}
\label{subsec:accuracy-to-solution}
%%%%%%%%%%%%%%%%%%%%%%%%%%%%%%%%%%%%%%%%%%%%%%%%%%%%%%%%%%%%

Performance comparisons should account for accuracy.

Suppose method \(A\) is twice as fast per iteration but requires ten times
as many iterations.

Then its total time may be worse.

A fair comparison holds the mathematical target approximately fixed:

\[
\boxed{
\text{same tolerance}
}
\]

or

\[
\boxed{
\text{same final error}.
}
\]

%%%%%%%%%%%%%%%%%%%%%%%%%%%%%%%%%%%%%%%%%%%%%%%%%%%%%%%%%%%%
\subsection{Performance Portability}
\label{subsec:hpmc-performance-portability}
%%%%%%%%%%%%%%%%%%%%%%%%%%%%%%%%%%%%%%%%%%%%%%%%%%%%%%%%%%%%

Performance portability means an algorithm remains reasonably efficient on
different hardware architectures.

This is challenging because CPUs, GPUs, and distributed systems favor
different execution patterns.

A portable numerical formulation should expose:

\[
\boxed{
\text{parallelism},
}
\]

\[
\boxed{
\text{locality},
}
\]

and

\[
\boxed{
\text{structured kernels}.
}
\]

%%%%%%%%%%%%%%%%%%%%%%%%%%%%%%%%%%%%%%%%%%%%%%%%%%%%%%%%%%%%
\subsection{Hardware-Aware Algorithms}
\label{subsec:hardware-aware-algorithms}
%%%%%%%%%%%%%%%%%%%%%%%%%%%%%%%%%%%%%%%%%%%%%%%%%%%%%%%%%%%%

Traditional algorithm analysis emphasizes arithmetic operation count.

Modern high-performance algorithms also consider:

\[
\boxed{
\text{memory traffic},
}
\]

\[
\boxed{
\text{communication},
}
\]

and

\[
\boxed{
\text{parallel depth}.
}
\]

Thus mathematically equivalent algorithms can have radically different
performance.

%%%%%%%%%%%%%%%%%%%%%%%%%%%%%%%%%%%%%%%%%%%%%%%%%%%%%%%%%%%%
\subsection{Algorithm Co-Design}
\label{subsec:algorithm-codesign}
%%%%%%%%%%%%%%%%%%%%%%%%%%%%%%%%%%%%%%%%%%%%%%%%%%%%%%%%%%%%

Algorithm co-design means developing mathematical algorithms with hardware
characteristics in mind.

Examples include:

\[
\boxed{
\text{communication-avoiding factorization},
}
\]

\[
\boxed{
\text{matrix-free PDE operators},
}
\]

and

\[
\boxed{
\text{mixed-precision iterative refinement}.
}
\]

The mathematical method and computational architecture are designed
together.

%%%%%%%%%%%%%%%%%%%%%%%%%%%%%%%%%%%%%%%%%%%%%%%%%%%%%%%%%%%%
\subsection{High-Performance Workflow}
\label{subsec:hpmc-workflow}
%%%%%%%%%%%%%%%%%%%%%%%%%%%%%%%%%%%%%%%%%%%%%%%%%%%%%%%%%%%%

A practical workflow is:

\begin{enumerate}
    \item formulate the mathematical problem;
    \item estimate its asymptotic computational complexity;
    \item identify dominant numerical kernels;
    \item determine memory requirements;
    \item determine expected communication requirements;
    \item choose sparse, dense, or matrix-free representation;
    \item begin with a verified serial or small-scale implementation;
    \item measure runtime and memory use;
    \item compute arithmetic intensity where useful;
    \item identify whether kernels are bandwidth or compute limited;
    \item improve data locality;
    \item vectorize regular arithmetic;
    \item introduce shared-memory parallelism;
    \item introduce distributed-memory decomposition when needed;
    \item minimize synchronization;
    \item reduce communication volume and message count;
    \item overlap communication with computation;
    \item evaluate accelerator suitability;
    \item minimize host-device data transfer;
    \item consider mixed precision when numerically safe;
    \item use scalable preconditioners;
    \item perform strong-scaling experiments;
    \item perform weak-scaling experiments;
    \item analyze load balance;
    \item profile communication and computation;
    \item implement checkpointing for long jobs;
    \item verify that parallelization preserves numerical correctness;
    \item document hardware, software, and numerical settings;
    \item report time to solution at the required accuracy.
\end{enumerate}

%%%%%%%%%%%%%%%%%%%%%%%%%%%%%%%%%%%%%%%%%%%%%%%%%%%%%%%%%%%%
\subsection{Choosing a High-Performance Strategy}
\label{subsec:choosing-hpmc-strategy}
%%%%%%%%%%%%%%%%%%%%%%%%%%%%%%%%%%%%%%%%%%%%%%%%%%%%%%%%%%%%

For dense matrix calculations, emphasize:

\[
\boxed{
\text{blocking}
+
\text{matrix-matrix kernels}.
}
\]

For sparse systems, emphasize:

\[
\boxed{
\text{memory efficiency}
+
\text{preconditioning}.
}
\]

For PDEs, emphasize:

\[
\boxed{
\text{domain decomposition}
+
\text{halo exchange}
+
\text{multigrid}.
}
\]

For Monte Carlo, emphasize:

\[
\boxed{
\text{independent ensemble parallelism}.
}
\]

For accelerators, emphasize:

\[
\boxed{
\text{large parallel kernels}
+
\text{data reuse}.
}
\]

%%%%%%%%%%%%%%%%%%%%%%%%%%%%%%%%%%%%%%%%%%%%%%%%%%%%%%%%%%%%
\subsection{Common Errors}
\label{subsec:hpmc-common-errors}
%%%%%%%%%%%%%%%%%%%%%%%%%%%%%%%%%%%%%%%%%%%%%%%%%%%%%%%%%%%%

\subsubsection{Optimizing Before Measuring}

Optimization should begin with profiling.

A visually complicated part of the code may contribute little to total
runtime.

%%%%%%%%%%%%%%%%%%%%%%%%%%%%%%%%%%%%%%%%%%%%%%%%%%%%%%%%%%%%
\subsubsection{Measuring FLOP/s Without Considering Time to Solution}

High arithmetic throughput does not guarantee faster solution of the
mathematical problem.

Always compare total runtime at comparable accuracy.

%%%%%%%%%%%%%%%%%%%%%%%%%%%%%%%%%%%%%%%%%%%%%%%%%%%%%%%%%%%%
\subsubsection{Ignoring Memory Bandwidth}

Many numerical kernels are limited by memory movement rather than floating-
point arithmetic.

Adding more arithmetic units may not improve them.

%%%%%%%%%%%%%%%%%%%%%%%%%%%%%%%%%%%%%%%%%%%%%%%%%%%%%%%%%%%%
\subsubsection{Ignoring Arithmetic Intensity}

An algorithm with low arithmetic intensity is unlikely to reach peak FLOP
performance.

Roofline-style reasoning can explain why.

%%%%%%%%%%%%%%%%%%%%%%%%%%%%%%%%%%%%%%%%%%%%%%%%%%%%%%%%%%%%
\subsubsection{Using Dense Storage for Sparse Problems}

A sparse matrix with \(O(N)\) nonzeros should not generally require
\(O(N^2)\) storage.

Representation choice determines scalability.

%%%%%%%%%%%%%%%%%%%%%%%%%%%%%%%%%%%%%%%%%%%%%%%%%%%%%%%%%%%%
\subsubsection{Storing a Matrix That Can Be Applied Matrix-Free}

For structured operators, explicit storage may waste memory and bandwidth.

Compute the operator action directly when appropriate.

%%%%%%%%%%%%%%%%%%%%%%%%%%%%%%%%%%%%%%%%%%%%%%%%%%%%%%%%%%%%
\subsubsection{Parallelizing Too Fine a Task}

If each task performs very little computation, scheduling and communication
overhead may dominate.

Increase granularity.

%%%%%%%%%%%%%%%%%%%%%%%%%%%%%%%%%%%%%%%%%%%%%%%%%%%%%%%%%%%%
\subsubsection{Assuming More Processors Always Reduce Runtime}

Strong scaling eventually saturates because of:

\[
\boxed{
\text{serial work},
}
\]

\[
\boxed{
\text{communication},
}
\]

and

\[
\boxed{
\text{synchronization}.
}
\]

%%%%%%%%%%%%%%%%%%%%%%%%%%%%%%%%%%%%%%%%%%%%%%%%%%%%%%%%%%%%
\subsubsection{Ignoring Load Balance}

A parallel program is limited by the slowest required worker.

Balanced workloads are essential.

%%%%%%%%%%%%%%%%%%%%%%%%%%%%%%%%%%%%%%%%%%%%%%%%%%%%%%%%%%%%
\subsubsection{Using Too Many Tiny Messages}

Repeated latency costs can dominate.

Batch communication when possible.

%%%%%%%%%%%%%%%%%%%%%%%%%%%%%%%%%%%%%%%%%%%%%%%%%%%%%%%%%%%%
\subsubsection{Using Global Reductions Excessively}

Global reductions synchronize the entire distributed system.

At large scale, they can dominate otherwise inexpensive iterative
operations.

%%%%%%%%%%%%%%%%%%%%%%%%%%%%%%%%%%%%%%%%%%%%%%%%%%%%%%%%%%%%
\subsubsection{Moving Data to a GPU for Tiny Calculations}

Transfer and launch overhead can exceed the saved computation time.

Accelerators require sufficient computational work per transfer.

%%%%%%%%%%%%%%%%%%%%%%%%%%%%%%%%%%%%%%%%%%%%%%%%%%%%%%%%%%%%
\subsubsection{Ignoring Kernel Fusion}

Writing intermediate arrays to memory can create unnecessary bandwidth
cost.

Fuse compatible operations when mathematically appropriate.

%%%%%%%%%%%%%%%%%%%%%%%%%%%%%%%%%%%%%%%%%%%%%%%%%%%%%%%%%%%%
\subsubsection{Using Low Precision Without Conditioning Analysis}

Mixed precision is powerful only when the underlying numerical problem can
tolerate lower precision.

Conditioning and convergence must be checked.

%%%%%%%%%%%%%%%%%%%%%%%%%%%%%%%%%%%%%%%%%%%%%%%%%%%%%%%%%%%%
\subsubsection{Using a Serial Preconditioner Inside a Parallel Solver}

The iterative method may be parallel while the preconditioner becomes the
scalability bottleneck.

Solver and preconditioner must be designed together.

%%%%%%%%%%%%%%%%%%%%%%%%%%%%%%%%%%%%%%%%%%%%%%%%%%%%%%%%%%%%
\subsubsection{Ignoring Coarse Multigrid Levels}

Fine-grid work may scale well while the coarse solve dominates at high
processor counts.

Coarse-level scalability must be analyzed separately.

%%%%%%%%%%%%%%%%%%%%%%%%%%%%%%%%%%%%%%%%%%%%%%%%%%%%%%%%%%%%
\subsubsection{Assuming an FFT Is Communication-Free}

Distributed FFTs can require expensive global data transpositions.

Arithmetic complexity alone does not capture parallel cost.

%%%%%%%%%%%%%%%%%%%%%%%%%%%%%%%%%%%%%%%%%%%%%%%%%%%%%%%%%%%%
\subsubsection{Using One Parallel Random Stream Everywhere}

Parallel stochastic computation requires separated or reproducibly
partitioned random-number streams.

%%%%%%%%%%%%%%%%%%%%%%%%%%%%%%%%%%%%%%%%%%%%%%%%%%%%%%%%%%%%
\subsubsection{Demanding Bitwise Reproducibility Without Need}

Exact bitwise reproducibility can impose additional performance costs.

Use it when scientifically necessary; otherwise define numerical tolerance.

%%%%%%%%%%%%%%%%%%%%%%%%%%%%%%%%%%%%%%%%%%%%%%%%%%%%%%%%%%%%
\subsubsection{Comparing Different Hardware Using Different Numerical Accuracy}

A faster result is not meaningful if it solves the problem less accurately.

Benchmark comparable mathematical outcomes.

%%%%%%%%%%%%%%%%%%%%%%%%%%%%%%%%%%%%%%%%%%%%%%%%%%%%%%%%%%%%
\subsubsection{Assuming Faster Hardware Fixes Poor Scaling}

Hardware cannot overcome an algorithm dominated by communication or serial
dependencies.

Redesign the algorithm when necessary.

%%%%%%%%%%%%%%%%%%%%%%%%%%%%%%%%%%%%%%%%%%%%%%%%%%%%%%%%%%%%
\subsection{Applications}
\label{subsec:hpmc-applications}
%%%%%%%%%%%%%%%%%%%%%%%%%%%%%%%%%%%%%%%%%%%%%%%%%%%%%%%%%%%%

\begin{applicationbox}

\textbf{Climate and Weather.}
Large atmosphere-ocean models solve coupled PDE systems over enormous
three-dimensional grids using distributed parallel computers.

\medskip

\textbf{Computational Fluid Dynamics.}
High-resolution flow simulation requires scalable sparse linear algebra,
domain decomposition, and accelerator computation.

\medskip

\textbf{Physics.}
Particle simulations, quantum systems, plasma models, and lattice
calculations frequently require massive parallelism.

\medskip

\textbf{Engineering.}
Structural mechanics, electromagnetics, heat transfer, and optimization
often involve millions of degrees of freedom.

\medskip

\textbf{Statistics and Data Analysis.}
Large least-squares, matrix factorization, Bayesian sampling, and
simulation problems benefit from distributed and accelerator-based
computation.

\medskip

\textbf{Optimization.}
Large-scale optimization uses parallel objective evaluations, distributed
gradients, matrix-free Hessian products, and scalable linear solvers.

\medskip

\textbf{Monte Carlo Simulation.}
Independent samples and trajectories provide natural coarse-grained
parallelism.

\medskip

\textbf{Mathematical Research.}
Large numerical experiments can explore conjectures, parameter spaces,
high-dimensional systems, and computational limits beyond manual
calculation.

\end{applicationbox}

%%%%%%%%%%%%%%%%%%%%%%%%%%%%%%%%%%%%%%%%%%%%%%%%%%%%%%%%%%%%
\subsection{Exercises}
\label{subsec:hpmc-exercises}
%%%%%%%%%%%%%%%%%%%%%%%%%%%%%%%%%%%%%%%%%%%%%%%%%%%%%%%%%%%%

\begin{exercise}[1]
Define high-performance mathematical computing.
\end{exercise}

\begin{exercise}[1]
Define arithmetic intensity.
\end{exercise}

\begin{exercise}[1]
Define memory bandwidth.
\end{exercise}

\begin{exercise}[1]
Define parallel speedup.
\end{exercise}

\begin{exercise}[1]
Define parallel efficiency.
\end{exercise}

\begin{exercise}[1]
Define strong scaling.
\end{exercise}

\begin{exercise}[1]
Define weak scaling.
\end{exercise}

\begin{exercise}[1]
Define communication latency.
\end{exercise}

\begin{exercise}[1]
Define matrix-free computation.
\end{exercise}

\begin{exercise}[1]
Define load imbalance.
\end{exercise}

\begin{exercise}[2]
Compute arithmetic intensity from a supplied FLOP count and byte count.
\end{exercise}

\begin{exercise}[2]
Use a roofline model to estimate an upper performance bound.
\end{exercise}

\begin{exercise}[2]
Determine whether a kernel is bandwidth bound or compute bound.
\end{exercise}

\begin{exercise}[3]
Compute the roofline crossover intensity

\[
I^*
=
\frac{
P_{\mathrm{peak}}
}{
B
}.
\]
\end{exercise}

\begin{exercise}[3]
Explain how cache blocking increases data reuse.
\end{exercise}

\begin{exercise}[3]
Explain why dense matrix multiplication has high arithmetic intensity.
\end{exercise}

\begin{exercise}[3]
Explain why sparse matrix-vector multiplication is often bandwidth bound.
\end{exercise}

\begin{exercise}[3]
Compare dense and sparse storage requirements.
\end{exercise}

\begin{exercise}[3]
Write a matrix-free implementation formula for a tridiagonal Laplacian.
\end{exercise}

\begin{exercise}[2]
Identify independent iterations in a vector update.
\end{exercise}

\begin{exercise}[3]
Explain how a race condition occurs during parallel summation.
\end{exercise}

\begin{exercise}[3]
Design a tree reduction for eight workers.
\end{exercise}

\begin{exercise}[2]
Compute parallel speedup from \(T_1\) and \(T_p\).
\end{exercise}

\begin{exercise}[2]
Compute parallel efficiency.
\end{exercise}

\begin{exercise}[3]
Apply Amdahl's law for a specified serial fraction.
\end{exercise}

\begin{exercise}[3]
Compute the maximum Amdahl speedup as

\[
p\to\infty.
\]
\end{exercise}

\begin{exercise}[3]
Explain why Amdahl's law is primarily a strong-scaling argument.
\end{exercise}

\begin{exercise}[3]
Explain Gustafson-type scaled-problem reasoning.
\end{exercise}

\begin{exercise}[3]
Design a strong-scaling experiment.
\end{exercise}

\begin{exercise}[3]
Design a weak-scaling experiment.
\end{exercise}

\begin{exercise}[3]
Explain why load imbalance reduces parallel efficiency.
\end{exercise}

\begin{exercise}[3]
Compare static and dynamic load balancing.
\end{exercise}

\begin{exercise}[2]
Use

\[
T_{\mathrm{comm}}
=
\alpha+\beta m
\]

to compare one large message with several small messages.
\end{exercise}

\begin{exercise}[3]
Explain why communication latency matters even when bandwidth is high.
\end{exercise}

\begin{exercise}[3]
Explain communication hiding.
\end{exercise}

\begin{exercise}[3]
Design a halo-exchange sequence for a one-dimensional domain
decomposition.
\end{exercise}

\begin{exercise}[3]
Extend the halo-exchange idea to a two-dimensional grid.
\end{exercise}

\begin{exercise}[3]
Explain the surface-to-volume effect.
\end{exercise}

\begin{exercise}[3]
Explain why distributed dot products require global reductions.
\end{exercise}

\begin{exercise}[3]
Explain why global reductions limit Krylov scaling.
\end{exercise}

\begin{exercise}[3]
Describe a pipelined Krylov strategy conceptually.
\end{exercise}

\begin{exercise}[3]
Explain why a preconditioner must also scale in parallel.
\end{exercise}

\begin{exercise}[3]
Explain the coarse-grid bottleneck in multigrid.
\end{exercise}

\begin{exercise}[3]
Explain why distributed FFTs require data redistribution.
\end{exercise}

\begin{exercise}[3]
Design a domain decomposition for a rectangular PDE grid.
\end{exercise}

\begin{exercise}[3]
Explain why unstructured meshes require graph partitioning.
\end{exercise}

\begin{exercise}[3]
Design a graph-partitioning objective balancing work and edge cuts.
\end{exercise}

\begin{exercise}[3]
Explain why adaptive mesh refinement creates dynamic load imbalance.
\end{exercise}

\begin{exercise}[2]
Design a parallel Monte Carlo estimator.
\end{exercise}

\begin{exercise}[3]
Explain how random streams should be separated across workers.
\end{exercise}

\begin{exercise}[3]
Design a parallel parameter sweep.
\end{exercise}

\begin{exercise}[3]
Describe nested parallelism for a simulation ensemble.
\end{exercise}

\begin{exercise}[3]
Parallelize the gradient

\[
\nabla F(x)
=
\sum_i
\nabla f_i(x).
\]
\end{exercise}

\begin{exercise}[3]
Explain distributed Hessian-vector products.
\end{exercise}

\begin{exercise}[3]
Construct a simple computational task graph.
\end{exercise}

\begin{exercise}[3]
Identify the critical path in a supplied task graph.
\end{exercise}

\begin{exercise}[3]
Compute theoretical parallelism from

\[
T_1/T_\infty.
\]
\end{exercise}

\begin{exercise}[3]
Explain asynchronous iteration.
\end{exercise}

\begin{exercise}[2]
Explain why accelerator data transfer can dominate a small computation.
\end{exercise}

\begin{exercise}[3]
Explain kernel-launch overhead.
\end{exercise}

\begin{exercise}[3]
Give an example of kernel fusion.
\end{exercise}

\begin{exercise}[3]
Explain why batching improves accelerator utilization.
\end{exercise}

\begin{exercise}[2]
Describe a mixed-precision algorithm.
\end{exercise}

\begin{exercise}[3]
Write the iterative-refinement correction equations.
\end{exercise}

\begin{exercise}[3]
Explain how conditioning limits mixed-precision accuracy.
\end{exercise}

\begin{exercise}[3]
Compare one-dimensional and two-dimensional distributed matrix layouts.
\end{exercise}

\begin{exercise}[3]
Explain why block matrix algorithms are useful in distributed dense linear
algebra.
\end{exercise}

\begin{exercise}[3]
Explain communication-avoiding QR conceptually.
\end{exercise}

\begin{exercise}[3]
Describe checkpoint/restart for a distributed simulation.
\end{exercise}

\begin{exercise}[3]
Explain the tradeoff in checkpoint frequency.
\end{exercise}

\begin{exercise}[3]
Explain resilience in large-scale computation.
\end{exercise}

\begin{exercise}[3]
Explain why floating-point parallel reductions may produce different last
digits.
\end{exercise}

\begin{exercise}[3]
Distinguish bitwise and numerical reproducibility.
\end{exercise}

\begin{exercise}[3]
Explain how a fixed reduction tree can improve reproducibility.
\end{exercise}

\begin{exercise}[3]
Design a profiling study for a distributed PDE solver.
\end{exercise}

\begin{exercise}[3]
Design a benchmark comparing CPU and GPU implementations at equal
accuracy.
\end{exercise}

\begin{exercise}[3]
Explain why time to solution is often more meaningful than raw FLOP/s.
\end{exercise}

\begin{exercise}[3]
Explain performance portability.
\end{exercise}

\begin{exercise}[3]
Give an example of hardware-aware numerical algorithm design.
\end{exercise}

\begin{exercise}[4]
Derive the basic roofline bound.
\end{exercise}

\begin{exercise}[4]
Estimate arithmetic intensity of blocked dense matrix multiplication.
\end{exercise}

\begin{exercise}[4]
Study cache-oblivious matrix algorithms.
\end{exercise}

\begin{exercise}[4]
Study SIMD vectorization of numerical kernels.
\end{exercise}

\begin{exercise}[4]
Analyze sparse matrix-vector arithmetic intensity.
\end{exercise}

\begin{exercise}[4]
Study communication lower bounds for dense linear algebra.
\end{exercise}

\begin{exercise}[4]
Study communication-avoiding LU and QR factorizations.
\end{exercise}

\begin{exercise}[4]
Analyze the communication complexity of parallel conjugate gradient.
\end{exercise}

\begin{exercise}[4]
Study pipelined Krylov methods.
\end{exercise}

\begin{exercise}[4]
Study scalable domain-decomposition preconditioners.
\end{exercise}

\begin{exercise}[4]
Study parallel geometric multigrid.
\end{exercise}

\begin{exercise}[4]
Study algebraic multigrid scalability.
\end{exercise}

\begin{exercise}[4]
Study distributed multidimensional FFT algorithms.
\end{exercise}

\begin{exercise}[4]
Analyze graph partitioning for unstructured finite-element meshes.
\end{exercise}

\begin{exercise}[4]
Study asynchronous iterative methods.
\end{exercise}

\begin{exercise}[4]
Study work-stealing task schedulers.
\end{exercise}

\begin{exercise}[4]
Analyze mixed-precision iterative refinement.
\end{exercise}

\begin{exercise}[4]
Study accelerator-aware sparse linear algebra.
\end{exercise}

\begin{exercise}[4]
Study matrix-free high-order finite-element methods.
\end{exercise}

\begin{exercise}[4]
Study scalable checkpointing algorithms.
\end{exercise}

\begin{exercise}[4]
Study reproducible parallel summation.
\end{exercise}

\begin{exercise}[4]
Analyze energy-to-solution as a performance metric.
\end{exercise}

\begin{exercise}[4]
Design a complete high-performance workflow for a large numerical PDE
problem.
\end{exercise}

%%%%%%%%%%%%%%%%%%%%%%%%%%%%%%%%%%%%%%%%%%%%%%%%%%%%%%%%%%%%
\subsection{Conceptual Summary}
\label{subsec:hpmc-summary}
%%%%%%%%%%%%%%%%%%%%%%%%%%%%%%%%%%%%%%%%%%%%%%%%%%%%%%%%%%%%

High-performance mathematical computing combines

\[
\boxed{
\text{numerical algorithms}
+
\text{parallelism}
+
\text{memory-aware computation}
+
\text{hardware architecture}.
}
\]

Performance is limited by both arithmetic and data movement.

The roofline model gives the conceptual bound

\[
\boxed{
P
\leq
\min
\left(
P_{\mathrm{peak}},
BI
\right).
}
\]

Parallel speedup is

\[
\boxed{
S_p
=
\frac{
T_1
}{
T_p
},
}
\]

and efficiency is

\[
\boxed{
E_p
=
\frac{
S_p
}{
p
}.
}
\]

Amdahl's law illustrates the effect of serial work:

\[
\boxed{
S_p
=
\frac{
1
}{
s+\frac{1-s}{p}
}.
}
\]

Scalable numerical algorithms minimize:

\[
\boxed{
\text{memory traffic},
}
\]

\[
\boxed{
\text{communication},
}
\]

and

\[
\boxed{
\text{synchronization}.
}
\]

Modern high-performance mathematics relies on:

\[
\boxed{
\text{blocking},
}
\]

\[
\boxed{
\text{vectorization},
}
\]

\[
\boxed{
\text{domain decomposition},
}
\]

\[
\boxed{
\text{accelerators},
}
\]

\[
\boxed{
\text{matrix-free methods},
}
\]

\[
\boxed{
\text{mixed precision},
}
\]

and

\[
\boxed{
\text{communication-avoiding algorithms}.
}
\]

The relevant goal is not peak machine speed alone, but

\[
\boxed{
\text{minimum reliable time to a mathematically accurate solution}.
}
\]

%%%%%%%%%%%%%%%%%%%%%%%%%%%%%%%%%%%%%%%%%%%%%%%%%%%%%%%%%%%%
\subsection{Section Connections}
\label{subsec:hpmc-connections}
%%%%%%%%%%%%%%%%%%%%%%%%%%%%%%%%%%%%%%%%%%%%%%%%%%%%%%%%%%%%

\chapterconnection{
High-Performance Mathematical Computing completes Chapter~11 by showing
how the numerical, symbolic, and simulation methods developed throughout
the chapter can be scaled to modern computational systems. Numerical
Linear Algebra supplies the dense, sparse, and iterative kernels;
Numerical ODEs and PDEs provide large time-dependent systems; Monte Carlo
Methods provide naturally parallel ensembles; Scientific Computing
provides the broader software and reproducibility framework; and
Mathematical Simulation provides application-scale computational
experiments. High-performance methods add communication-aware algorithms,
distributed data structures, accelerators, matrix-free operators, mixed
precision, scalable solvers, and performance modeling. These computational
tools support the formal foundations developed in the next chapter,
Mathematical Logic and Foundations, while also providing the computational
infrastructure for the applied modeling topics developed later in the
text.
}

%%%%%%%%%%%%%%%%%%%%%%%%%%%%%%%%%%%%%%%%%%%%%%%%%%%%%%%%%%%%
% SECTION SOURCE TRACKING
%%%%%%%%%%%%%%%%%%%%%%%%%%%%%%%%%%%%%%%%%%%%%%%%%%%%%%%%%%%%

%------------------------------------------------------------
% 11.11 High-Performance Mathematical Computing
%------------------------------------------------------------
%
% Source basis:
%   - Standard high-performance computing.
%   - Standard parallel numerical algorithms.
%   - Standard communication-aware and accelerator-oriented scientific
%     computing principles.
%
% Used for:
%   - high-performance mathematical computing;
%   - performance layers;
%   - algorithmic improvement versus hardware speed;
%   - FLOP counts and throughput;
%   - peak and sustained performance;
%   - memory bandwidth;
%   - arithmetic intensity;
%   - roofline modeling;
%   - bandwidth- and compute-bound regimes;
%   - cache blocking and tiling;
%   - SIMD and vectorization;
%   - dense matrix multiplication;
%   - matrix-vector performance;
%   - sparse matrix-vector multiplication;
%   - sparse storage;
%   - structured sparse operators;
%   - matrix-free computation;
%   - shared-memory parallelism;
%   - parallel loop independence;
%   - race conditions;
%   - reductions and synchronization;
%   - distributed-memory parallelism;
%   - point-to-point and collective communication;
%   - communication cost modeling;
%   - latency and bandwidth;
%   - communication avoidance;
%   - communication hiding;
%   - nonblocking communication;
%   - domain decomposition;
%   - halo exchange;
%   - surface-to-volume scaling;
%   - speedup and efficiency;
%   - Amdahl's law;
%   - strong and weak scaling;
%   - Gustafson-type scaling;
%   - load balancing;
%   - accelerator computing;
%   - GPU execution;
%   - host-device transfer;
%   - kernel launch overhead;
%   - kernel fusion;
%   - batching;
%   - mixed precision;
%   - iterative refinement;
%   - distributed dense matrices;
%   - block matrix distributions;
%   - parallel matrix multiplication;
%   - parallel LU, QR, and Cholesky;
%   - communication-avoiding factorization;
%   - parallel Krylov methods;
%   - global reductions;
%   - pipelined Krylov methods;
%   - scalable preconditioning;
%   - parallel multigrid;
%   - coarse-grid bottlenecks;
%   - parallel FFTs;
%   - distributed PDE solvers;
%   - structured and unstructured grids;
%   - graph partitioning;
%   - adaptive mesh load balancing;
%   - parallel Monte Carlo;
%   - random streams;
%   - ensemble and nested parallelism;
%   - parallel optimization;
%   - distributed gradients and Hessian-vector products;
%   - task-based parallelism;
%   - dependency graphs;
%   - critical paths;
%   - work and span;
%   - asynchronous iteration;
%   - fault tolerance;
%   - checkpointing and resilience;
%   - floating-point reproducibility;
%   - profiling and scaling experiments;
%   - energy efficiency;
%   - time to solution;
%   - performance portability;
%   - hardware-aware algorithm design;
%   - applications and exercises.
%
% Notes:
%   - This completes Chapter 11.
%   - Chapter 12 begins with Formal Logic.
%   - Scientific Computing appears in Section 11.7.
%   - Numerical PDEs appear in Section 11.5.
%   - Monte Carlo Methods appear in Section 11.6.
%
%%%%%%%%%%%%%%%%%%%%%%%%%%%%%%%%%%%%%%%%%%%%%%%%%%%%%%%%%%%%